% Modernised reading — fonds Alexandre Grothendieck, shelfmark 79, pages 1-296
% (the whole folder).
%
% This is an INTERPRETATION, not a transcription. It re-reads the pages in
% today's notation and vocabulary, states the hypotheses the manuscript leaves
% implicit, and drops the critical apparatus so that the mathematics reads
% straight through. Everything the brackets and underlines carried moves into
% a footnote. It is held to being mathematically correct as it stands; where
% the manuscript is loose or mistaken, the true statement is what appears here
% and the footnote says what the page has. In any dispute of substance the
% transcription governs: whoever cites Grothendieck cites batch-01.fr.tex to
% batch-15.fr.tex (pp. 1-20, 21-40, ..., 261-280, 281-296).
%
% Pass: Opus 5.5 (claude-opus-5-5), 2026-10-11 — first pass, unchecked against the pages by a human.
% The folder was read whole, all fifteen batches, batch by batch with notes;
% this is one document for the shelfmark. It was made on Opus 5.5 and has not
% been measured against the reference reading of folder 115 (made on Opus 5).
% The computations the reading asserts beyond the page were checked by hand
% while writing, not by machine: the matrices of pp. 27 and 32 and their
% products; the counts N'/n, N'/2, N'/3 and the genus for n = 3 to 7; the
% normaliser computations of pp. 93-95; the rank of Z cap G_0^+ (p. 151);
% the comparison of the abstract constructions of pp. 102-269 with
% GL(2, Z/n)/{±1} (torus, Borel, big Bruhat cell, det as the "cyclotomic
% character"); the counterexamples given in footnotes (an additive map of
% F_4^2, the pair (2, 0) over Q, the orbits for n = 6, the triangle over Z/5,
% the regular map {3,6}_(3,0)). The literature named in footnotes (Farey
% tessellation, modular curves X(n), Klein's map, Serre's Arbres, Wilson's
% operators and map exponents, Sabidussi coset graphs, Tits systems, the
% generation of SL_2 of a semi-local ring by elementary matrices) is cited
% from memory and was not checked against the sources.
%
% Scope. All fifteen batches are transcribed. Not transcribed, and not read
% here: page 1 is a cover whose words title the folder; pages 271, 274 and
% 281 are blank; pages 287, 289, 291 and 296 are typed university notices
% with nothing in his hand. Page 273 carries one formula. Pages 121-125 (the
% end of § VII), the long diagonal marginal notes throughout, and the scratch
% sheets 282-295 are read only as far as the transcription supports, which is
% partly; the reading says so where it leans on them.
%
% Spine. One object, the triangulated surface drawn by a free module M of rank
% 2 over Lambda = Z or Z/n, and one question, how much of M the surface knows:
%   pp. 3-20    § I: triangles x+y+z = 0, the polyhedra P_M and SP_M (gauge
%               e), stars as affine lines, connectedness iff Lambda = Lambda_0;
%   pp. 21-26   the theorem SL(M)' = Aut^+(SP_M), GL^±(M)' = Aut(SP_M); the
%               counts and the genus (0 iff n <= 5, 1 iff n = 6);
%   pp. 26-37   § II: the cartographic group and PSL(2, Z), PGL(2, Z); the
%               flag automorphisms as matrices; the « formulaire 1981 »;
%               universality of SP_M over Z;
%   pp. 37-60   § III: no gauge; the disjoint union of the phi(n)/2 polyhedra
%               (bi-icosahedron for n = 5), strict automorphisms, the
%               commutant, two attempts to pass from one structure to all;
%   pp. 61-70   § IV: arbitrary Lambda, types, cases 2 = 0 and 2 != 0;
%   pp. 71-91   § V: « Récapitulation »: P~_M on M^*, the lemma and theorems,
%               Lambda^* acting, the quotient P_M, groupoid equivalences, a
%               programme a)-f);
%   pp. 92-101  § VI: the commutant of SL(M) and SL(M)'; the octahedral
%               exception 2Lambda = {0,2}; Z'_Lambda read on the stars;
%   pp. 102-128 § VII: « Structures triangulaires congruentielles »: axioms
%               T1-T3, T_0 1-T_0 6, the twisted graphs A_lambda, To 7;
%   pp. 129-158 § VIII: « Axiomatique : reformulation »: To 1-To 6, the
%               « fait mirifique » (Z/n)^*/±1 acting on n-gons, the cyclotomic
%               character, the splitting of a vertex stabiliser;
%   pp. 159-177 § IX: « Graphes à homothéties » and their group theory;
%   pp. 178-221 § X: antipodisms; invariant contours; graphs as
%               vertex-star incidences with a « jauge »; prégraphes;
%   pp. 222-270 § XI: « Récapitulation »: prégraphes, rays, the conditions of
%               isomorphy and homogeneity, the group-theoretic data
%               (H, U_0, T_H, sigma), N_H; triangulating graphs (begun);
%   pp. 272-295 loose sheets: a table for n = 2-5; « Détermination d'un
%               sommet par son étoile »; pencil and scratch computations.
%
% Notation fixed for the whole document.
%  - Lambda a nonzero commutative ring, Lambda_0 the image of Z in it, M free
%    of rank 2. M^* is the set of elements of M that are part of a basis (the
%    page's M^*, S~_M, M^{**}: NOT a dual); (det M)^* the set of generators of
%    det M = Lambda^2 M, the « jauges »; the dual module is M^v.
%  - S_M = M^*/{±1}; F_M the faces {x., y., z.} with x+y+z = 0 and (x, y) a
%    basis, A_M the edges; Pi_M = (det M)^*/{±1} the « types »; P_{M,i} the
%    polyhedron of type i (the page's SP_M for i = {±e}); P_M the union of all
%    types on the one vertex set S_M; P~_M the same on M^* before dividing by
%    ±1. The § III disjoint union is written P_M^sqcup (the page reuses P_M,
%    and its margin says so).
%  - GL, SL modern for his Gl, Sl; GL^±(M) the det = ±1 part; a prime means
%    « modulo {±1} »: SL(M)' = SL(M)/{±1}. Z'_Lambda = Lambda^*/{±1}, Z'_n for
%    Lambda = Z/n.
%  - tau_0, tau_1, tau_infty change the vertex, edge, face of a flag; rho_0 =
%    tau_infty tau_1, rho_1 = tau_0 tau_infty, rho_infty = tau_1 tau_0 the
%    rotations about a vertex, an edge, a face (pp. 27, 112, 164).
%  - In the abstract sections, S, A_0, F_0, G_0 = Aut, G_0^+, Z (the group of
%    « homothéties »), A_lambda = {{s, lambda t} : {s, t} in A_0}, pi the set
%    of structures, G, delta : G -> Z, Et(s) the star of s. The names of the
%    axioms (T_0 i, To i, T i, GH 1) are the page's, section by section.
%
% Substantive departures from the pages, each footnoted where it occurs:
%  p. 8    the union in (19) runs over F_M, not P_3(S_M).
%  p. 16   the components of a star are the cosets y + Lambda_0 x, their set
%          is L_x/Lambda_0 x (page: « L_x wedge Lambda/Lambda_0 »).
%  p. 21   connectedness of SP_M for Lambda = Z (left « sûrement vrai ») is
%          supplied; the theorem of p. 23 needs it.
%  p. 24   oriented edges of SP_M = Repsp(M)', not Rep(M)'.
%  p. 36   algebraic proofs that the cartographic group is PSL(2, Z) exist.
%  p. 40   S_{M,i} is in bijection with M^*/{±1}, not M^*.
%  p. 44   (110) e): one endpoint moved, else the type moves by lambda^2.
%  p. 46   pi is the image of (det M)^* in (det M)/{±1}; it is not « the
%          unique orbit of cardinal phi(n)/2 » when n = 2 mod 4, n >= 6.
%  p. 63   Ker delta_M = GL^±(M)/{±1}, which contains SL(M)/{±1} with index 2.
%  p. 66   « même sur Z » is wrong: over Z (and Z/3, Z/4, Z/6) the triangles
%          of A_M are faces; the relation fails as soon as Lambda^* != {±1}.
%  p. 68   the guess about characteristic 2 is refuted by p. 74, (15).
%  p. 76   the Lemma needs (a, b) unimodular in Lambda_0, not in Lambda.
%  p. 83   an additive automorphism of M need not preserve bases (F_4);
%          the true statement is for semilinear automorphisms.
%  p. 87   the components of a star have card Lambda_0, not card Lambda.
%  p. 99   the map Z -> Z'_Lambda is T_mu |-> mu, no exponent.
%  p. 101  the torsor is under Lambda^*, not Lambda.
%  p. 142  Z'_9 has order 3 (page: 4); Z'_15 order 4 (page: 12, doubtful).
%  p. 154  Z = (Z/n)^*/{±1}; it is T that is isomorphic to (Z/n)^*.
%  p. 175  in (57 a) the relation of order 3 is (tau_1 tau_0)^3.
%  p. 182  « ordres pairs » must read « finis » (the case (16 b)).
%  p. 242  (70 b) must read u(s_0, t_0) = (lambda s_0, lambda^-1 t_0).
%  p. 249  T_0, T_1 are defined the other way round from Ker delta_0, delta_1.
%  p. 252  « T'_0 » is T_1.
%  p. 275  « convexe » must read « connexe ».
%
% Modern names given and footnoted as ours: Farey tessellation, principal
% congruence subgroup, modular curve X(n), regular map {3, n}, Klein's map of
% genus 3, hemi-icosahedron, regular map {3,6}_(3,0), Wilson's operators and
% map exponents, polarity without absolute points, coset graph, arc-regular
% graph, voltage graph (regular covering of a graph), Tits system (BN-pair),
% Borel subgroup, torus, big Bruhat cell, Weil pairing. « Groupe
% cartographique » is his own term.
%
% Announced and never established: connectedness of SP_M over Z (p. 21,
% supplied here); an algebraic proof of p. 36; the universality statement of
% p. 37 (breaks off); an a priori reason for (101) => (95) (p. 43); the
% argument « à expliciter » of p. 46; recovering the map wedge from the other
% data (p. 50); the question of p. 44; the exactness of (24), p. 82 (supplied
% for Lambda = Lambda_0); the programme of pp. 90-91; the checks « à
% vérifier » of pp. 98-99; whether T_0 4 follows from the other axioms
% (p. 114); the case n even of p. 120 (answered pp. 149-150); the question
% for n = 7 (p. 142); the abelianisation of GL'(2, Z/8)^± (p. 147); the
% translation of GH 1 (p. 177); condition b) of p. 189; the torsor claim of
% p. 212 (withdrawn by him); the formula B_H cap N_H = T_H (p. 269); § XI.4
% (p. 270, begun only).
\documentclass[11pt,a4paper]{article}
\input{../preamble/grothendieck}

\folder{79}
\pages{1}{296}
\dating{[à partir de 1980]}
\watermark{Édition de démonstration\\interprétation personnelle de l'œuvre}
\foldertitle{2-polyèdres réguliers triangulés, et modules libres de rangs 2 sur les anneaux $\Lambda=\mathbb{Z}/n\mathbb{Z}$ et $\Lambda=\mathbb{Z}$ : notes manuscrites (s.d.) — lecture modernisée du dossier entier}

\begin{document}

\section*{Résumé}

\begin{resume}
Ces feuillets partent d'un objet très simple, le plan des points à
coordonnées entières, et lui font dessiner une surface. On prend les vecteurs
$(a, b)$ à coordonnées premières entre elles, en identifiant chacun à son
opposé ; on relie deux d'entre eux quand ils forment une base du réseau
(déterminant $\pm 1$) ; trois vecteurs dont la somme est nulle forment un
triangle. Ce qu'on obtient est une triangulation infinie, où chaque sommet est
entouré d'une infinité de triangles. Si l'on remplace les entiers par
l'arithmétique « de l'horloge », les entiers modulo $n$, la même recette donne
une surface finie : un tétraèdre pour $n = 3$, un octaèdre pour $n = 4$, un
icosaèdre pour $n = 5$, un tore à douze sommets pour $n = 6$, puis des surfaces
à plusieurs trous. Le dossier part de là : un module de rang $2$ sur
$\mathbf{Z}$ ou sur $\mathbf{Z}/n\mathbf{Z}$ dessine de lui-même un polyèdre
régulier aux faces triangulaires.

Le premier résultat dit que ce dessin est fidèle : les symétries du polyèdre
sont exactement les matrices $2 \times 2$ de déterminant $\pm 1$, au signe près. La
question qui occupe la suite est la question inverse, et elle est plus
difficile : que faut-il savoir du polyèdre pour retrouver le module ? Quand on
oublie le choix d'une unité de surface (le « déterminant » qui vaut $1$),
plusieurs polyèdres s'installent sur les mêmes sommets --- deux icosaèdres
entrelacés pour $n = 5$ --- et l'on veut comprendre comment l'un détermine les
autres, et comment la multiplication par un nombre inversible modulo $n$ se
lit sur le dessin.

L'idée qui arrive est d'une grande simplicité. Autour de chaque sommet, les
voisins forment un polygone à $n$ côtés. Mais un $n$-gone peut se parcourir de
plusieurs façons : on peut sauter d'un sommet au suivant, ou de deux en deux,
comme on trace un pentagramme sur les sommets d'un pentagone. Les sauts
permis, comptés au signe près, forment le groupe $(\mathbf{Z}/n\mathbf{Z})^{*}
/ \{\pm 1\}$, qui agit ainsi sur les polygones de tout ensemble à $n$ éléments
--- le « fait mirifique » de la page 138. Multiplier les vecteurs par un
inversible $\lambda$ revient à retracer chaque étoile en sautant de $\lambda$
en $\lambda$ : l'arithmétique de l'anneau se lit sur les étoiles. Il en sort
une suite de tentatives d'axiomatique --- quelles structures combinatoires
proviennent d'un module ? --- reprise quatre fois, avec un vocabulaire qui
s'affine : graphes « à homothéties », dualité entre sommets et étoiles,
« prégraphes » et « jauges », torseurs. Elle aboutit à de la théorie des
groupes pure, où reparaissent, sans être nommés, les sous-groupes qui
charpentent les matrices $2 \times 2$ : le tore des matrices diagonales, les
matrices triangulaires, et la décomposition de presque toute matrice en un
produit triangulaire supérieur, diagonal, triangulaire inférieur.

On lit aussi le dossier pour sa manière de travailler. Il garde ses erreurs
et ses repentirs (« relation fausse !! », un « faux » posé sur un énoncé de la
veille, un « Ouf » après une vérification) ; il fait « bon poids
d'hypothèses, quitte à vérifier par la suite » ; il s'étonne de retrouver par
pure combinatoire des formules obtenues l'année précédente « par voie
topologico-transcendante », et de ne pas connaître de démonstration purement
algébrique d'une identité qu'il juge remarquable. Il ne conclut pas : il
s'arrête sur des feuilles de calcul, au moment où il voit que la
multiplication par $\lambda$ ne se lit sur une étoile qu'au carré près, mais
se lit exactement sur les incidences entre deux étoiles.

Les noms à chercher ensuite : pavage de Farey et groupe modulaire, sous-groupes
de congruence et courbes modulaires $X(n)$, cartes régulières (et la carte de
Klein), groupe cartographique, exposants d'une carte, graphes de classes et
graphes arc-transitifs, systèmes de Tits, caractère cyclotomique.

\keywords{unimodular vector, free module of rank 2, Farey tessellation, regular map, triangulated surface, modular group, extended modular group, cartographic group, principal congruence subgroup, modular curve X(n), Platonic solids, Klein quartic, centralizer of a permutation group, star polygon, map exponent, Wilson operator, hemi-icosahedron, bipyramid, polarity, coset graph, arc-transitive graph, voltage graph, torsor, projective line over a ring, cyclotomic character, Borel subgroup, Bruhat decomposition, BN-pair, elementary matrices}
\end{resume}

\pagerange{1}{296}
\section*{Le fil du dossier, et les conventions}

La couverture (page 1) porte le titre que reprend l'inventaire : « 2-polyèdres
réguliers triangulés, et modules libres de rang 2 sur les anneaux $\Lambda =
\mathbf{Z}/n\mathbf{Z}$ et $\Lambda = \mathbf{Z}$ ». La page 2 est une table des
matières de sa main, qui renvoie à sa propre pagination (page 1 de l'auteur =
page 3 du fonds, jusqu'à sa page 138 = page 270)\footnote{À partir de la
ligne « Axiomatique des structures triangulées congruentielles (1) », la table
décale d'un cran les numéros romains par rapport aux titres du corps, et
plusieurs de ses numéros de page ne coïncident pas avec ceux des titres ; les
lignes VII à X y sont très surchargées. On suit ici les titres du corps.}. Le
texte est continu, numéroté par paragraphes de I à XI, chaque paragraphe
repartant souvent de sa propre numérotation des formules ; on cite les
formules par leur numéro et leur page. Les pages 272 à 295 sont des feuilles
détachées.

La datation de l'inventaire est « à partir de 1980 ». Deux indices internes
la précisent sans la fixer : la page 33 compare les calculs à ceux d'un
« formulaire 1981 », et la page 35 dit « l'an dernier j'ai choisi
délibérément » une convention de ce formulaire, ce qui place au moins ces
pages en 1982 ou après\footnote{Inférence de notre part ; le formulaire n'est
pas dans ce dossier. Un mot souligné en tête de la page 240, lu
« vendredi », est peut-être une date ; la lecture est très incertaine.}.

\subsection*{Les stations}

\begin{itemize}
  \item \emph{Pages 3 à 26 : § I et le théorème.} Pour $M$ libre de rang $2$
    sur $\Lambda$, les triangles $\{x, y, z\}$, $x + y + z = 0$, de vecteurs
    deux à deux en base ; le polyèdre $P_M$ des sommets $M^{*}/\{\pm 1\}$ et,
    une fois choisie une jauge $e$, le polyèdre $SP_M$ ; l'étoile d'un sommet
    est une droite affine sur $\Lambda$, connexe exactement quand $\Lambda =
    \mathbf{Z}$ ou $\mathbf{Z}/n\mathbf{Z}$ ; orientation ; théorème
    $\mathrm{SL}(M)' \simeq \mathrm{Aut}^{+}(SP_M)$ ; nombres de sommets,
    arêtes, faces et genre (pages 21 à 26).
  \item \emph{Pages 26 à 37 : § II.} Le groupe cartographique triangulé, son
    isomorphisme avec $\mathrm{PSL}(2, \mathbf{Z})$ et $\mathrm{PGL}(2,
    \mathbf{Z})$, les automorphismes de drapeaux comme matrices, le
    « formulaire 1981 », l'universalité de $SP_M$ sur $\mathbf{Z}$.
  \item \emph{Pages 37 à 60 : § III.} Sans jauge : la réunion disjointe des
    $\varphi(n)/2$ polyèdres, ses automorphismes stricts, le torseur des types
    et l'application $\wedge$, le commutant, deux tentatives pour déduire
    toutes les structures d'une seule.
  \item \emph{Pages 61 à 70 : § IV.} Anneau quelconque : types, cas $2 = 0$
    et $2 \neq 0$, ce que les arêtes déterminent.
  \item \emph{Pages 71 à 91 : § V, « Récapitulation ».} La structure
    $\widetilde{P}_M$ sur $M^{*}$, un lemme, deux théorèmes, l'action de
    $\Lambda^{*}$, le quotient $P_M$, des équivalences de groupoïdes, un
    programme en six points.
  \item \emph{Pages 92 à 101 : § VI.} Le commutant de $\mathrm{SL}(M)$ et de
    $\mathrm{SL}(M)'$ ; l'exception « octaédrique » ; la lecture de
    $Z'_\Lambda$ sur les étoiles.
  \item \emph{Pages 102 à 158 : §§ VII et VIII.} Deux axiomatiques
    successives des structures triangulées « congruentielles », puis
    « prégroupales » ; le « fait mirifique » ; le caractère cyclotomique ; le
    scindage du stabilisateur d'un sommet.
  \item \emph{Pages 159 à 221 : §§ IX et X.} Graphes à homothéties et leur
    traduction en théorie des groupes ; antipodismes ; contours invariants ;
    graphes vus comme incidences entre sommets et étoiles ; prégraphes et
    jauges.
  \item \emph{Pages 222 à 270 : § XI, « Récapitulation ».} Prégraphes,
    rayons, conditions d'isomorphie et d'homogénéité, et la description d'un
    multigraphe homogène par un groupe $H$ et trois sous-groupes ; les graphes
    triangulants, à peine commencés.
  \item \emph{Pages 272 à 295 : feuilles détachées.} Un tableau pour $n = 2$
    à $5$ ; « Détermination d'un sommet par son étoile » ; brouillons.
\end{itemize}

Ce ne sont pas des chapitres d'un exposé mais des reprises. Le § V récapitule
les §§ I à IV dans un autre langage ; les §§ VII, VIII, IX et XI sont quatre
\emph{moutures} de la même axiomatique, chacune avec ses axiomes et ses
noms ; on les garde distinctes, et on renvoie de l'une à l'autre là où elles se
recouvrent.

\subsection*{Conventions}

$\Lambda$ est un anneau commutatif non nul, $\Lambda_0$ l'image de $\mathbf{Z}$
dans $\Lambda$ ; on dit avec la page que $\Lambda$ est « absolu » quand
$\Lambda = \Lambda_0$, c'est-à-dire $\Lambda \simeq \mathbf{Z}$ ou
$\mathbf{Z}/n\mathbf{Z}$. $M$ est un $\Lambda$-module libre de rang $2$,
$\mathrm{Rep}(M)$ l'ensemble de ses bases ordonnées (« repères »). On note,
comme la page, $M^{*}$ l'ensemble des éléments de $M$ qui font partie d'une
base, et $(\det M)^{*}$ l'ensemble des générateurs du module inversible $\det M
= \Lambda^2 M$ --- les « jauges »\footnote{L'astérisque désigne donc ici les
éléments « inversibles », non un dual ; le dual de $M$ est noté $M^{\vee}$.
La page 6 définit d'abord $M^{*}$ par la non-annulation en tout point de
$\operatorname{Spec}\Lambda$, et note que la définition par les bases lui
équivaut sous $\operatorname{Pic}(\Lambda) = 0$ ; à partir de la page 9 (et
de la page 61) c'est la définition par les bases qui est retenue, et c'est
celle qu'on suit.}. Pour $x \in M$, $\dot{x}$ est sa classe modulo $\pm 1$.

$S_M = M^{*}/\{\pm 1\}$ ; une \emph{face} est une partie $\{\dot x, \dot y,
\dot z\}$ avec $x + y + z = 0$ et $(x, y)$ base, une \emph{arête} une partie
$\{\dot x, \dot y\}$ avec $(x, y)$ base ; $F_M$, $A_M$ leurs ensembles.
Le \emph{type} d'une arête ou d'une face est la classe $\pm(x \wedge y) \in
\Pi_M = (\det M)^{*}/\{\pm 1\}$. On note $P_{M,i}$ le polyèdre formé des faces
de type $i$ (c'est le $SP_M$ du § I pour $i = \{\pm e\}$), $P_M$ la réunion de
tous les types sur l'ensemble de sommets $S_M$, et $\widetilde{P}_M$ la même
construction sur $M^{*}$ avant division par $\pm 1$. La réunion
\emph{disjointe} des $P_{M,i}$ que considère le § III est notée
$P_M^{\sqcup}$\footnote{La page 37 réutilise le nom $P_M$ pour cette réunion
disjointe, et sa marge note qu'elle a « changé la notation p.r. à (23) ».}.
$Z'_\Lambda = \Lambda^{*}/\{\pm 1\}$, et $Z'_n$ pour $\Lambda =
\mathbf{Z}/n\mathbf{Z}$ ; c'est un groupe d'ordre $\varphi(n)/2$ pour $n \geq
3$.

Pour les groupes, on écrit $\mathrm{GL}$, $\mathrm{SL}$ là où la page écrit
$\mathrm{Gl}$, $\mathrm{Sl}$ ; $\mathrm{GL}^{\pm}(M)$ est le sous-groupe des
automorphismes de déterminant $\pm 1$, et un prime signifie « modulo $\{\pm
1\}$ » : $\mathrm{SL}(M)' = \mathrm{SL}(M)/\{\pm 1\}$. Pour une carte
triangulée, $\tau_0$, $\tau_1$, $\tau_\infty$ changent respectivement le
sommet, l'arête, la face d'un drapeau, et $\rho_0 = \tau_\infty\tau_1$,
$\rho_1 = \tau_0\tau_\infty$, $\rho_\infty = \tau_1\tau_0$ sont les rotations
autour d'un sommet, d'un milieu d'arête, d'un centre de face (pages 27, 112,
164).

Dans les §§ VII à XI, plus abstraits, $S$ est un ensemble de sommets, $A_0
\subset \mathfrak{P}_2(S)$ un ensemble d'arêtes, $F_0 \subset
\mathfrak{P}_3(S)$ un ensemble de faces, $G_0$ le groupe des automorphismes,
$G_0^{+}$ celui des automorphismes directs, $\mathrm{Et}(s)$ l'étoile
(l'ensemble des voisins) de $s$. Le groupe noté $Z$ y joue le rôle de
$Z'_\Lambda$ : ses éléments sont des permutations de $S$, les
« homothéties », qui ne respectent pas $A_0$ mais le tordent en
\[
  A_\lambda = \bigl\{ \{s, \lambda t\} \bigm| \{s, t\} \in A_0 \bigr\}
  \qquad (\lambda \in Z) .
\]
C'est la construction centrale de tout le dossier, et elle a un modèle simple
qu'il faut avoir en tête : sur $S_M$, la multiplication $T_\mu$ par un scalaire
$\mu \in \Lambda^{*}$, appliquée à \emph{une seule} extrémité d'une arête,
change son type $i$ en $\mu i$, tandis qu'appliquée aux deux extrémités elle le
change en $\mu^2 i$. Toute la difficulté du dossier se loge dans cette
différence entre $\mu$ et $\mu^2$. Les lettres $\sigma$ et $\lambda$ servent à
plusieurs usages d'un paragraphe à l'autre (pour $\lambda$ : les homothéties,
et au § II les matrices $\lambda_0$, $\lambda_1$, $\lambda_\infty$) ; on les
redéfinit sur place.

\pagerange{3}{10}
\section*{I. Les triangles d'un module de rang 2 (pages 3 à 20)}

\subsection*{Triangles, repères, orientation (pages 3 à 6)}

Soient $x, y, z \in M$ avec $x + y + z = 0$. Alors $(x, y)$ est une base si et
seulement si $(y, z)$ en est une, si et seulement si $(z, x)$ en est une. Soit
$B(M) \subset M^3$ l'ensemble de ces triplets ; le groupe symétrique
$\mathfrak{S}_3$ y opère librement, et le quotient $\widetilde{F}(M) =
B(M)/\mathfrak{S}_3$ est l'ensemble des parties $\{x, y, z\}$ de somme nulle
dont deux éléments quelconques forment une base. L'application $(x, y, z)
\mapsto (x, y)$ identifie $B(M)$ à $\mathrm{Rep}(M) =
\mathrm{Isom}_\Lambda(\Lambda^2, M)$ ; l'action de $\mathfrak{S}_3$ s'y lit
comme l'action naturelle sur $V'_\Lambda = \Lambda^3/\Lambda$ (quotient par la
diagonale), via $B(M) \simeq \mathrm{Isom}_\Lambda(V'_\Lambda, M)$ et
l'isomorphisme $\Lambda^2 \simeq \Lambda^3/\Lambda$, $(x, y) \mapsto (x, y, 0)$.

Pour $(x, y, z) \in B(M)$, on a $x \wedge y = y \wedge z = z \wedge x \in (\det
M)^{*}$. Cet élément ne dépend que de l'ordre circulaire $\omega$ défini sur $A
= \{x, y, z\}$ par l'ordre choisi, et change de signe avec lui : on a une
application $\alpha_A$ de l'ensemble $\underline{\omega}_A$ des deux ordres
circulaires dans $(\det M)^{*}$, avec $\alpha_A(\omega^{-1}) =
-\alpha_A(\omega)$. Si l'on fixe une jauge $e \in (\det M)^{*}$ et que $2 \neq
0$ dans $\Lambda$, les triangles $A$ pour lesquels $\alpha_A$ prend les
valeurs $\pm e$ reçoivent une orientation canonique $\omega_A$, celle pour
laquelle $\alpha_A(\omega_A) = e$.

Tout élément de $M^{*}$ appartient à un tel triangle : si $(x, y)$ est une
base, $\{x, y, -(x+y)\}$ en est un ; et deux éléments de $M^{*}$ appartiennent à
au plus un triangle\footnote{Page 6. La page démontre l'appartenance sous
l'hypothèse $\operatorname{Pic}(\Lambda) = 0$ (par exemple $\Lambda$
semi-local), avec sa première définition de $M^{*}$ ; avec la définition par
les bases, elle est immédiate.}.

\pagerange{7}{10}
\subsection*{Le polyèdre $P_M$ (pages 7 à 10)}

On suppose désormais $2 \cdot 1_\Lambda \neq 0$. L'involution $-\mathrm{id}_M$
opère alors sans point fixe sur $M^{*}$, et l'on pose $S_M = M^{*}/\{\pm 1\}$,
$F_M = \widetilde{F}_M/\{\pm 1\}$. L'application $F_M \to \mathfrak{P}_3(S_M)$
est injective : si $\{x, y, z\}$ et $\{x, y, -z\}$ étaient tous deux des
triangles, on aurait $2z = 0$, donc $2 = 0$ puisque $z$ fait partie d'une base.
On identifie donc $F_M$ à un ensemble de parties à trois éléments de $S_M$, de
réunion $S_M$\footnote{La page 8 écrit la réunion sur $A \in
\mathfrak{P}_3(S_M)$ ; elle doit porter sur les faces $A \in F_M$, comme le dit
sa propre note marginale.}, et l'on note $A_M$ l'ensemble des paires $\{\dot
x, \dot y\}$ avec $(x, y)$ base.

\emph{Proposition} (pages 9 et 10). Toute arête $\{\dot x, \dot y\}$ est
contenue dans exactement deux faces, $\{\dot x, \dot y, (x+y)^{\cdot}\}$ et
$\{\dot x, \dot y, (x-y)^{\cdot}\}$ : les triangles qui la relèvent sont $\{\varepsilon
x, \varepsilon' y, -(\varepsilon x + \varepsilon' y)\}$, $\varepsilon,
\varepsilon' = \pm 1$, à $(\varepsilon, \varepsilon') \mapsto (-\varepsilon,
-\varepsilon')$ près. Ainsi $P_M = (S_M, A_M, F_M)$ est un « 2-polyèdre
combinatoire triangulé » au sens faible : toute arête borde exactement deux
faces\footnote{Une note marginale de la page 10 avertit qu'il manque la
condition de connexité (du contour d'un sommet) ; elle est étudiée aux pages
14 à 17.}. Si $\Lambda$ est semi-local, deux sommets sont toujours à distance
au plus $2$, et $P_M$ est connexe.

\pagerange{11}{13}
\subsection*{La jauge et le polyèdre $SP_M$ (pages 11 à 13)}

Fixons une jauge $e$. Les faces de type $\{\pm e\}$ --- les $\{\dot x, \dot y,
\dot z\}$ avec $x \wedge y = \pm e$ --- forment $SF_M$, leurs arêtes $SA_M$. Si
$f$, $f'$ sont les deux faces qui contiennent l'arête $a = \{\dot x, \dot y\}$,
les applications $a \to (\det M)^{*}$ qu'elles induisent par $\alpha_f$,
$\alpha_{f'}$ diffèrent par le signe (la première envoie $\dot x$ sur $x \wedge
y$, la seconde sur $x \wedge (-y)$). Il en résulte que $a$, $f$ et $f'$ ont le
même type, et que
\[
  SP_M = (S_M, SA_M, SF_M)
\]
est encore un polyèdre triangulé dont toute arête borde deux faces ; c'est le
$P_{M, \{\pm e\}}$ de nos conventions.

\pagerange{14}{17}
\subsection*{L'étoile d'un sommet (pages 14 à 17)}

Soit $s = \dot x$. Les voisins de $s$ dans $SP_M$ correspondent aux $y$ tels
que $x \wedge y = e$ ; ils forment la droite affine
\[
  L_x = \{ y \in M \mid x \wedge y = e \},
\]
espace homogène sous les translations par $\Lambda x$. Deux voisins $\dot y$,
$\dot z$ ($y, z \in L_x$) forment une face avec $s$ si et seulement si $y - z =
\pm x$ : de $x + \varepsilon' y + \varepsilon'' z = 0$ et $x \wedge (y - z) =
0$, l'égalité $\varepsilon' = \varepsilon''$ donnerait $2e = 0$. L'étoile de
$s$ est donc la droite $L_x$ munie de la structure de « contour » où chaque
point est relié à son translaté par $\pm x$, et ses composantes connexes sont
les classes $y + \Lambda_0 x$\footnote{La page 16 écrit l'ensemble de ces
composantes « $L_x \wedge_\Lambda \Lambda/\Lambda_0$ ». C'est le quotient
$L_x/\Lambda_0 x$, torseur sous le groupe additif $\Lambda/\Lambda_0$ (comme
$x$ fait partie d'une base, $\Lambda x \simeq \Lambda$).}. D'où :
\begin{quote}
  l'étoile de chaque sommet de $SP_M$ est connexe si et seulement si $\Lambda =
  \Lambda_0$, c'est-à-dire $\Lambda \simeq \mathbf{Z}$ ou $\Lambda \simeq
  \mathbf{Z}/n\mathbf{Z}$ avec $n \geq 3$.
\end{quote}
Pour $\Lambda = \mathbf{Z}/n\mathbf{Z}$, l'étoile est un $n$-gone ; pour
$\Lambda = \mathbf{Z}$, c'est une ligne infinie de voisins. Le polyèdre $P_M$
sans jauge, réunion des polyèdres de tous les types, n'a d'étoiles connexes que
s'il n'y a qu'un type, c'est-à-dire $\Lambda^{*} = \{\pm 1\}$ : les seuls cas
sont $\Lambda = \mathbf{Z}$ et $\mathbf{Z}/n\mathbf{Z}$, $n \in \{3, 4, 6\}$
(page 17).

\pagerange{17}{20}
\subsection*{Orientation et connexité (pages 17 à 20)}

$SP_M$ ne change pas si l'on remplace $e$ par $-e$ : il est stable par
$\mathrm{GL}^{\pm}(M)$. Les orientations $\omega_f$ des faces (celles où
$\alpha_f(\omega_f) = e$) sont compatibles le long des arêtes, par
l'anticommutativité vue plus haut, et définissent une orientation canonique
$\mu$ de $SP_M$, qui dépend de $e$ et que $g \in \mathrm{GL}^{\pm}(M)$ envoie
sur $\det(g)\,\mu$.

Pour la connexité, il s'agit de relier deux éléments de $M^{*}$ par une chaîne
$z_0, \dots, z_N$ avec $z_i \wedge z_{i+1} = \pm e$. Si $(x, y)$ est une base
avec $e = \lambda\, x \wedge y$, $\lambda \in \Lambda^{*}$, l'élément $z =
\lambda x + \lambda y$ vérifie $x \wedge z = z \wedge y = e$ : il suffit donc de
relier par des chaînes de bases, ce qui est toujours possible sur un anneau
semi-local. On trouve même que deux sommets de $SP_M$ sont à distance au plus
$3$ quand $\Lambda$ est semi-local, la borne étant atteinte pour l'icosaèdre
($\Lambda = \mathbf{Z}/5\mathbf{Z}$). Pour $\Lambda = \mathbf{Z}$, la page 21
juge « sûrement vrai » que $SP_M$ est connexe, sans s'y attarder\footnote{C'est
vrai, et c'est l'algorithme d'Euclide : un vecteur primitif $(a, b)$ de
$\mathbf{Z}^2$ distinct de $\pm(1, 0), \pm(0, 1)$ est la somme de deux vecteurs
primitifs plus courts formant une base de déterminant $\pm 1$ (ses
« parents » de Farey), ce qui le relie par récurrence à $(1, 0)$. L'argument
est de nous ; le théorème de la page 23 en a besoin pour $\Lambda =
\mathbf{Z}$.}.

\pagerange{21}{26}
\section*{Le théorème d'automorphismes et le dénombrement (pages 21 à 26)}

Un automorphisme $u$ de $M$ qui fixe les trois sommets d'une face est $\pm
\mathrm{id}$ : quitte à changer son signe, $ux = x$, $uy = \varepsilon' y$, $uz
= \varepsilon'' z$, et $x = -y - z = -\varepsilon' y - \varepsilon'' z$ force
$\varepsilon' = \varepsilon'' = 1$. Les arêtes orientées de $SP_M$
correspondent aux bases symplectiques $(x, y)$, $x \wedge y = e$, modulo $\pm
1$\footnote{La page 24, (58), écrit $\overrightarrow{SA}_M \simeq
\mathrm{Rep}(M)'$ ; il faut lire $\mathrm{Repsp}(M, e)' = \{(x, y) \mid x
\wedge y = e\}/\{\pm 1\}$, comme en (51), page 22.}, sur lesquelles
$\mathrm{SL}(M)'$ opère simplement transitivement. Comme un automorphisme d'un
2-polyèdre connexe au sens strict est connu par son effet sur un drapeau, on
obtient :

\emph{Théorème} (page 23). Si $\Lambda = \mathbf{Z}$ ou
$\mathbf{Z}/n\mathbf{Z}$ ($n \geq 3$),
\[
  \mathrm{SL}(M)' \xrightarrow{\ \sim\ } \mathrm{Aut}^{+}(SP_M), \qquad
  \mathrm{GL}^{\pm}(M)' \xrightarrow{\ \sim\ } \mathrm{Aut}(SP_M),
\]
et $SP_M$ est un 2-polyèdre triangulé orienté connexe, \emph{régulier} au sens
orienté comme au sens non orienté\footnote{Ce sont des noms classiques, que la
page ne donne pas : pour $\Lambda = \mathbf{Z}$, $SP_M$ est le \emph{pavage de
Farey} du plan hyperbolique (sommets $\mathbf{P}^1(\mathbf{Q})$, arêtes les
paires $a/c$, $b/d$ avec $ad - bc = \pm 1$) ; pour $\Lambda =
\mathbf{Z}/n\mathbf{Z}$, c'est son quotient par le sous-groupe de congruence
principal $\Gamma(n)$, la carte régulière de type $\{3, n\}$ tracée sur la
courbe modulaire $X(n)$, connue depuis Klein. Cités de mémoire.}.

Pour $\Lambda = \mathbf{Z}/n\mathbf{Z}$, $n = \prod p^{\nu_p}$, les cardinaux se
calculent par le lemme chinois :
\[
  N = \operatorname{card}\mathrm{SL}(2, \mathbf{Z}/n\mathbf{Z}) = \prod_p
  p(p^2 - 1)\,p^{3(\nu_p - 1)} = n^3 \prod_{p \mid n}\Bigl(1 -
  \frac{1}{p^2}\Bigr), \qquad N' = N/2 ,
\]
$\operatorname{card} M^{*} = N/n$, et $SP_M$ a $N'/n$ sommets, $N'/2$ arêtes,
$N'/3$ faces ; tous les sommets sont d'ordre $n$, toutes les faces d'ordre $3$.
Sa caractéristique d'Euler est
\[
  \chi = N'\Bigl(\frac{1}{n} - \frac{1}{2} + \frac{1}{3}\Bigr) = N'\Bigl(\frac{1}{n}
  - \frac{1}{6}\Bigr) = 2 - 2g , \qquad
  g = 1 + \frac{N'}{2}\Bigl(\frac{1}{6} - \frac{1}{n}\Bigr) ,
\]
d'où $g = 0$ si et seulement si $n \leq 5$, $g = 1$ si et seulement si $n = 6$,
$g \geq 2$ si et seulement si $n \geq 7$. Concrètement\footnote{Le tableau est
de nous, vérifié à la main ; la page s'arrête à la trichotomie. La formule est
celle, classique, du genre de $X(n)$ pour $n \geq 3$ ; pour $n = 7$ c'est la
carte de Klein, tracée sur la quartique de Klein. Le dossier 70 (pages 61 et
62) dresse le tableau complémentaire : le quotient du groupe triangulaire par
la puissance $n$-ième de la rotation autour d'un sommet est
$\mathrm{SL}(2, \mathbf{Z}/n)/\{\pm 1\}$ exactement pour $n \leq 5$.} :
\[
\begin{array}{c|cccc|l}
  n & N' & \text{sommets} & \text{arêtes} & \text{faces} & \\ \hline
  3 & 12 & 4 & 6 & 4 & \text{tétraèdre} \\
  4 & 24 & 6 & 12 & 8 & \text{octaèdre} \\
  5 & 60 & 12 & 30 & 20 & \text{icosaèdre} \\
  6 & 72 & 12 & 36 & 24 & \text{tore } (g = 1) \\
  7 & 168 & 24 & 84 & 56 & g = 3
\end{array}
\]
La page annonce qu'elle reviendra sur les cas $n \leq 6$, et passe d'abord à
la fonctorialité.

\pagerange{26}{37}
\section*{II. Le groupe cartographique et $\mathrm{SL}(2, \mathbf{Z})$ (pages 26 à 37)}

Les cartes triangulées orientées se décrivent comme des ensembles sous le
\emph{groupe cartographique} orienté triangulé
\[
  \mathfrak{G}^{\circ}_{3,\infty} = \langle \rho_0, \rho_1, \rho_\infty \mid
  \rho_1^2 = \rho_\infty^3 = 1,\ \rho_\infty\rho_1\rho_0 = 1 \rangle ,
\]
et les cartes non orientées comme des ensembles sous
\[
  \mathfrak{G}_{3,\infty} = \langle \tau_0, \tau_1, \tau_\infty \mid \tau_0^2 =
  \tau_1^2 = \tau_\infty^2 = (\tau_0\tau_\infty)^2 = (\tau_0\tau_1)^3 = 1
  \rangle ,
\]
avec $\rho_\infty = \tau_1\tau_0$, $\rho_0 = \tau_\infty\tau_1$, $\rho_1 =
\tau_0\tau_\infty$. On a des isomorphismes $\mathfrak{G}^{\circ}_{3,\infty}
\simeq \mathrm{PSL}(2, \mathbf{Z})$ et $\mathfrak{G}_{3,\infty} \simeq
\mathrm{PGL}(2, \mathbf{Z})$, donnés « suivant un formulaire » par
\[
  \lambda_1 = \begin{pmatrix} 0 & 1 \\ -1 & 0 \end{pmatrix}, \quad
  \lambda_\infty = \begin{pmatrix} 1 & -1 \\ 1 & 0 \end{pmatrix}, \quad
  \lambda_0 = \begin{pmatrix} 1 & -1 \\ 0 & 1 \end{pmatrix} ;
\]
\[
  r_\infty = \begin{pmatrix} -1 & 0 \\ 0 & 1 \end{pmatrix}, \quad
  r_0 = \begin{pmatrix} 0 & 1 \\ 1 & 0 \end{pmatrix}, \quad
  r_1 = \begin{pmatrix} -1 & 1 \\ 0 & 1 \end{pmatrix} ,
\]
images de $\rho_1, \rho_\infty, \rho_0$ et de $\tau_\infty, \tau_0, \tau_1$.
On vérifie $\lambda_0 = r_\infty r_1$, $\lambda_1 = r_0 r_\infty$,
$\lambda_\infty = r_1 r_0$, et $\lambda_\infty\lambda_1\lambda_0 = 1$ dans
$\mathrm{SL}(2, \mathbf{Z})$ même\footnote{Vérifications de nous. Les lettres
inscrites à côté des matrices (« $\sigma$ », « $\rho^{-1}$ », « $\varepsilon_0$ »)
sont lues sous réserve.}.

\pagerange{28}{30}
\subsection*{Le polyèdre comme ensemble sous le groupe cartographique (pages 28 à 30)}

$\mathrm{Rep}(M)' = \mathrm{Isom}(\Lambda^2, M)/\{\pm 1\}$ est un torseur à
droite sous $\mathrm{GL}(2, \Lambda)'$, donc, par $\mathbf{Z} \to \Lambda$, un
ensemble sous $\mathrm{GL}(2, \mathbf{Z})' \simeq \mathfrak{G}_{3,\infty}$ :
une structure triangulée, et de même une structure orientée sur les repères
symplectiques modulo $\pm 1$, ensemble sous $\mathrm{SL}(2, \mathbf{Z})'$. Pour
$\Lambda = \mathbf{Z}$ ou $\mathbf{Z}/n\mathbf{Z}$, on retrouve « sauf erreur »
le polyèdre $SP_M$ ; comme $\mathrm{SL}(2, \mathbf{Z}) \to \mathrm{SL}(2,
\Lambda)$ est alors surjectif, la structure est connexe et régulière ; dans le
cas général, ses composantes sont des polyèdres réguliers isomorphes, de
groupe d'automorphismes directs $\mathrm{SL}(2, \Lambda_0)'$ (« à
vérifier »)\footnote{La surjectivité de $\mathrm{SL}(2, \mathbf{Z}) \to
\mathrm{SL}(2, \mathbf{Z}/n\mathbf{Z})$ est classique ; elle est utilisée sans
démonstration.}.

\pagerange{31}{36}
\subsection*{Les automorphismes de drapeaux comme matrices (pages 31 à 36)}

Un repère direct $r = (s, a, f)$ de $SP_M$ revient à une base $(x, y)$ avec $x
\wedge y = e$, modulo signe : $s = \dot x$, $a = \{\dot x, \dot y\}$, $f =
\{\dot x, \dot y, (x+y)^{\cdot}\}$ ; la seconde face de $a$ a pour troisième
sommet $(x - y)^{\cdot}$. Les automorphismes $\tau_0^r$, $\tau_1^r$,
$\tau_\infty^r$ qui échangent $r$ avec ses drapeaux voisins sont induits par
les automorphismes antisymplectiques de matrices, dans la base $(x, y)$,
\[
  \tau_0^r : \begin{pmatrix} 0 & 1 \\ 1 & 0 \end{pmatrix} = r_0, \qquad
  \tau_1^r : \begin{pmatrix} 1 & -1 \\ 0 & -1 \end{pmatrix} = -r_1, \qquad
  \tau_\infty^r : \begin{pmatrix} 1 & 0 \\ 0 & -1 \end{pmatrix} = -r_\infty
\]
($\tau_0^r$ échange $x$ et $y$ ; $\tau_1^r$ fixe $x$ et échange $y$ et $z =
-x-y$ ; $\tau_\infty^r$ fixe $x$ et change $y$ en $-y$), d'où $\rho_0^r \mapsto
\lambda_0$, $\rho_1^r \mapsto -\lambda_1$, $\rho_\infty^r \mapsto
-\lambda_\infty$. Aux signes près, sans importance dans $\mathrm{SL}(2,
\Lambda)'$, on « tombe pile sur les formules du formulaire 1981, qui avaient
été obtenues par voie topologico-transcendante (classique) relativement
délicate » : cela confirme que le choix des générateurs $\rho =
\lambda_\infty^{-1}$ et $\sigma = \lambda_1$ de $\mathrm{SL}(2, \mathbf{Z})$,
présenté par $\rho^3 = \sigma^2 = \rho^{-3}$, était « le plus naturel
possible »\footnote{Cette présentation dit $\rho^6 = 1$ et $\sigma^2 = \rho^3$
(élément central $-1$) : c'est la présentation classique de $\mathrm{SL}(2,
\mathbf{Z})$ comme produit amalgamé $\mathbf{Z}/4 *_{\mathbf{Z}/2}
\mathbf{Z}/6$. Le dernier membre est lu sous réserve. Le formulaire de 1981
n'est pas dans ce dossier.}.

Relevant ensuite les homomorphismes à $\mathrm{SL}(2, \mathbf{Z})$ et
$\mathrm{GL}(2, \mathbf{Z})$ depuis le groupe libre $\Pi_{0,3} = \langle
\rho_0, \rho_1, \rho_\infty \mid \rho_\infty\rho_1\rho_0 = 1\rangle$ et le
produit libre $\Pi^2_{0,3}$ de trois groupes d'ordre $2$, la page montre qu'une
fois imposé que l'image de $\rho_0$ soit unipotente, il reste deux choix :
celui du formulaire (87), où $\lambda_\infty$ est d'ordre $6$, et celui (88)
que la géométrie vient de donner, où $-\lambda_\infty$ est d'ordre $3$.
« L'an dernier j'ai choisi délibérément » le premier, « pour avoir un élément
d'ordre $6$, moins spécial ».

\pagerange{36}{37}
\subsection*{Une identité remarquable, et l'universalité (pages 36 et 37)}

La construction du § I fournit donc $r_0, r_1, r_\infty \in \mathrm{GL}(2,
\mathbf{Z})$ et $\lambda_0, \lambda_1, \lambda_\infty \in \mathrm{SL}(2,
\mathbf{Z})$ vérifiant les relations des groupes cartographiques, et
$\lambda_\infty^3 = \lambda_1^2 = -1$. Que les homomorphismes
$\mathfrak{G}^{\circ}_{3,\infty} \to \mathrm{SL}(2, \mathbf{Z})'$ et
$\mathfrak{G}_{3,\infty} \to \mathrm{GL}(2, \mathbf{Z})'$ soient des
isomorphismes est, dit la page, une « identité remarquable, dont je ne connais
pas de démonstration par voie purement algébrique »\footnote{Il en existe :
par l'action de $\mathrm{PSL}(2, \mathbf{Z})$ sur un arbre (Serre,
\emph{Arbres, amalgames, $\mathrm{SL}_2$}, 1977), ou par l'algorithme d'Euclide
et un argument de forme normale. Cité de mémoire ; rien n'indique que la page
ait pu les avoir en vue.}. Elle l'énonce comme une propriété
d'\emph{universalité} : si l'on choisit un repère direct $r_0$ de $SP_M$ pour
$M = \mathbf{Z}^2$, alors pour toute structure 2-polyédrale triangulée
orientée $P$ et tout repère $r$ de $P$, il existe un homomorphisme et un seul
$SP_M \to P$ envoyant $r_0$ sur $r$. La phrase s'interrompt sur « tel
que »\footnote{L'énoncé, de nous : le pavage de Farey est la carte triangulée
universelle, et toute carte triangulée en est le quotient par un sous-groupe
de $\mathrm{PSL}(2, \mathbf{Z})$, un « homomorphisme » de cartes étant un
revêtement. C'est la correspondance entre sous-groupes du groupe modulaire et
cartes triangulées, voisine des dessins d'enfants (voir les dossiers 81 et
88).}.

\pagerange{37}{43}
\section*{III. Sans jauge : la famille des polyèdres (pages 37 à 60)}

\subsection*{Les structures supplémentaires et le théorème (pages 37 à 43)}

Soit $M$ libre de rang $2$ sur $\Lambda \simeq \mathbf{Z}$ ou
$\mathbf{Z}/n\mathbf{Z}$, $n \geq 3$, sans jauge choisie. Les structures
« semi-symplectiques » sont les éléments de $\Pi_M = (\det M)^{*}/\{\pm 1\}$ ;
à chacune $i$ correspond le polyèdre régulier $P_{M,i}$, dont l'ensemble des
orientations s'identifie à $i$ lui-même (une paire $\{e, -e\}$). Sur
$\mathbf{Z}$ il n'y en a qu'un. Sur $\mathbf{Z}/n\mathbf{Z}$ il y en a
$\varphi(n)/2$, et leur réunion disjointe $P_M^{\sqcup}$ est connexe si et
seulement si $\varphi(n) = 2$, $n \in \{3, 4, 6\}$ ; le premier cas non connexe
est $n = 5$, où l'on trouve un « bi-icosaèdre ».

Pour $\Lambda$ fini et $P_M^{\sqcup}$ non connexe, la page cherche des
structures supplémentaires dont le groupe d'automorphismes serait
$\mathrm{GL}(M)'$ tout entier. Elle en retient trois :
\begin{itemize}
  \item (95) l'action de $Z'_\Lambda = \Lambda^{*}/\{\pm 1\}$, sous-groupe
    central de $\mathrm{GL}(M)'$ (« égal au centre sauf erreur »), par les
    homothéties $T_\lambda$ ;
  \item (97) la structure de torseur sous $Z'_\Lambda$ de l'ensemble
    $\pi_0(P_M^{\sqcup}) \simeq \Pi_M$ des composantes ; on a $T_g\, i =
    \det(g) \cdot i$ pour $g \in \mathrm{GL}(M)'$, donc $T_\lambda\, i =
    \lambda^2 i$ pour $\lambda \in Z'_\Lambda$ ;
  \item (101) le système transitif de bijections $\kappa_{j,i}$ entre les
    ensembles de sommets des composantes, tous identifiés à $S_M$\footnote{La
    page 40, (100), écrit $S_{M,i} \simeq M^{*}$ ; l'ensemble des sommets
    d'une composante est $M^{*}/\{\pm 1\} = S_M$.}.
\end{itemize}
Un automorphisme \emph{strict} est un automorphisme de $P_M^{\sqcup}$ qui les
respecte. Le groupe $\mathrm{Autstr}(P_M^{\sqcup})$ s'envoie sur $Z'_\Lambda$
par son action sur les composantes, et l'on a un homomorphisme de suites
exactes
\begin{tikzcd}[column sep=small]
1 \arrow[r] & \mathrm{GL}^{\pm}(M)' \arrow[r] \arrow[d] & \mathrm{GL}(M)' \arrow[r] \arrow[d] & \Lambda^{*}/\{\pm 1\} \arrow[r] \arrow[d] & 1 \\
1 \arrow[r] & \mathrm{Autstr}^{0}(P_M^{\sqcup}) \arrow[r] & \mathrm{Autstr}(P_M^{\sqcup}) \arrow[r] & Z'_\Lambda \arrow[r] & 1
\end{tikzcd}
\emph{Théorème} (page 42 ; la démonstration se lit dans l'ordre des pages 42,
41, 43). C'est un isomorphisme. L'injectivité vient du théorème de la page
23. Pour la surjectivité, un automorphisme strict qui fixe chaque composante
induit sur l'une d'elles un $T_g$, $g \in \mathrm{GL}^{\pm}(M)'$ ; une fois
corrigé par $g$, il est l'identité sur une composante, donc sur toutes par
(101).

La démonstration n'a utilisé que (97) et (101), de sorte que (101) entraîne
(95) a posteriori : « il serait bon d'en comprendre la raison a priori ». Le
§ V montrera pourquoi le problème disparaît quand on met tous les polyèdres sur
le \emph{même} ensemble de sommets.

\pagerange{43}{46}
\subsection*{Le type d'objet obtenu (pages 43 à 46)}

La page récapitule la structure en données (105) : un ensemble $\mathbb{S}$
(la valeur commune des ensembles de sommets), un torseur $\pi$ sous
$Z'_\Lambda$, et une application $\Sigma$ de $\pi$ dans l'ensemble des
structures triangulées sur $\mathbb{S}$, avec les propriétés : a) chaque
$\Sigma(i)$ est un 2-polyèdre triangulé connexe ; b) leurs groupes
d'automorphismes $G$ sont égaux, et $G$, $G^{+}$ sont simplement transitifs
sur les drapeaux (resp. les drapeaux directs) ; c) $i \neq j$ entraîne que
$\Sigma(i)$, $\Sigma(j)$ n'ont aucune arête ni aucune face commune ; d) le
groupe $A$ des automorphismes de toute la donnée se surjecte sur $Z'_\Lambda$ ;
e) pour $\lambda \in Z'_\Lambda$ il existe une unique permutation $T_\lambda$
de $\mathbb{S}$ telle que $\{s, t\}$ soit une arête de $\Sigma(i)$ si et
seulement si $\{s, T_\lambda t\}$ est une arête de $\Sigma(\lambda
i)$\footnote{La page 44 écrit la condition e) en appliquant $T_\lambda$ aux
\emph{deux} extrémités ; ainsi écrite, elle relie $\Sigma(i)$ à
$\Sigma(\lambda^2 i)$, par (99), et non à $\Sigma(\lambda i)$. La démonstration
d'unicité qui suit (page 45) ne déplace qu'une extrémité, comme la formule
(125) de la page 49 : c'est la forme qu'on retient.}. Une note marginale
demande s'il existe d'autres structures triangulées sur $\mathbb{S}$, isomorphes
à celles-là et de mêmes groupes $G$, $G^{+}$.

L'unicité de $T_\lambda$ repose sur ce qu'un sommet est déterminé par ses
voisinages : si deux éléments $u, u'$ de $M^{*}$ ont les mêmes voisins pour
chaque type, alors $u' = \pm u$, ce qu'on ramène (« à expliciter ») au fait que
sur un corps ayant au moins trois éléments le complémentaire de deux droites
engendre le plan\footnote{L'énoncé exact, avec son exception, est le lemme des
pages 276 et 277 : il est vrai sauf quand $2\Lambda = \{0, 2\} \neq \{0\}$,
c'est-à-dire, pour $\Lambda = \mathbf{Z}/n\mathbf{Z}$, sauf pour $n = 4$, où
deux sommets opposés de l'octaèdre ont la même étoile.}.

\pagerange{46}{50}
\subsection*{Le torseur des types et l'application $\wedge$ (pages 46 à 50)}

À l'aide de $\pi = \Pi_M$ on reconstitue la « droite modulo $\pm 1$ » $L' =
(\det M)/\{\pm 1\}$, sur laquelle $Z'_n$ opère, et $\pi$ s'y plonge comme
l'image de $(\det M)^{*}$\footnote{La page 46 caractérise $\pi$ comme
« l'unique orbite de cardinal $\varphi(n)/2$ » de $Z'_n$ dans $L'$. C'est faux
quand $n \equiv 2 \pmod 4$, $n \geq 6$ : l'orbite de $2$ a aussi
$\varphi(n/2)/2 = \varphi(n)/2$ éléments ; pour $n = 6$, les quatre orbites
$\{0\}$, $\{\pm 1\}$, $\{\pm 2\}$, $\{3\}$ sont toutes de cardinal $1 =
\varphi(6)/2$. Ce qui est vrai, et suffit à la suite, c'est que $\pi$ est
l'orbite des générateurs.}. Sur $S = M^{*}/\{\pm 1\}$, l'application symétrique
$\wedge : S \times S \to L'$, induite par $x \wedge y$, détermine pour chaque $i
\in \pi$ la structure $\Sigma(i)$ : les arêtes sont les $\{x, y\}$ avec $x \wedge
y = i$, les faces les triplets dont toutes les paires sont des arêtes de type
$i$\footnote{Que ces triangles soient bien des faces (c'est-à-dire $z = \pm x
\pm y$) est démontré plus loin, page 67, (178).}. Chaque arête de type $i$ est
dans exactement deux faces, et l'on conjecture que l'étoile d'un sommet est un
$n$-gone.

Pour définir les $\Sigma(i)$, il suffit de connaître la partie $A$ de
$\mathfrak{P}_2(S)$ formée des arêtes de tous types et l'application « type »
$t : A \to \pi$, dont les fibres sont les $A_{\Sigma(i)}$. La structure de
torseur sur $\pi$ correspond alors à une action $\lambda \mapsto T_\lambda$ de
$Z'_n$ sur $S$ telle que
\[
  s \wedge T_\lambda t = \lambda\,(s \wedge t) ,
\]
uniquement déterminée puisqu'un sommet est connu par son étoile ; appliquée aux
deux extrémités, $T_\lambda$ multiplie le type par $\lambda^2$. Une note
marginale ajoute que les deux données --- l'action de $Z'_n$ sur $\pi$ et
l'action $T$ sur $S$ --- se déterminent l'une l'autre. La question de
reconstruire $\wedge$ elle-même à partir des autres données est esquissée puis
barrée : « il faut y revenir ultérieurement ».

\pagerange{50}{53}
\subsection*{Une structure détermine-t-elle les autres ? (pages 50 à 53)}

Soit $P_0 = (S, A_0, F_0)$ un 2-polyèdre triangulé connexe régulier fini,
$G_0 = \mathrm{Aut}(P_0) \subset \mathfrak{S}_S$, et $Z$ le commutant de $G_0$
dans $\mathfrak{S}_S$ ; pour $P_0 = P_{M,i}$, ce commutant se réduirait à
$Z'_n$ (« à vérifier »)\footnote{Le § VI (pages 92 à 95) le démontre pour le
commutant de $\mathrm{SL}(M)'$, sauf pour $n = 4$ ; pour $n = 4$ l'élément
supplémentaire est l'antipodie de l'octaèdre, qui appartient à $G_0$ et
disparaît dans le quotient $Z/\mathfrak{z}$ introduit ici.}. On pose
$\mathfrak{z} = Z \cap G_0 = \mathrm{Cent}(G_0)$ et $Z' = Z/\mathfrak{z}$.
Comme $G_0$ est transitif sur les arêtes et sur les faces, une structure
$\lambda\Sigma_0$ transformée par $\lambda \in Z$ est soit égale à $\Sigma_0$,
soit sans arête ni face commune avec elle. Deux structures régulières
$\Sigma$, $\Sigma'$ sur $S$ sont dites \emph{équivalentes} si $\Sigma' =
\lambda\Sigma$ pour un $\lambda$ commutant à $\mathrm{Aut}(S, \Sigma)$ ; une
classe d'équivalence $\Pi$ est un torseur sous $Z'$, et l'on a une suite exacte
$1 \to G \to \mathrm{Aut}(S, \Pi) \to Z' \to 1$.

Mais cette analyse, appliquée à $P_{M,i}$, ne donne que les $P_{M,\lambda^2 i}$,
puisque $T_\lambda$, appliqué aux deux extrémités, multiplie le type par
$\lambda^2$ : « on est encore loin du compte ! »\footnote{Pour $n = 5$ par
exemple, les carrés de $Z'_5 = \mathbf{F}_5^{*}/\{\pm 1\}$ sont triviaux, et
l'on n'atteint pas le second icosaèdre du bi-icosaèdre.}

\pagerange{54}{60}
\subsection*{Seconde tentative : la torsion d'une seule extrémité (pages 54 à 60)}

La page introduit un groupe a priori plus petit, $Z_0$ : celui des $\lambda$
qui centralisent $G_0$ et tels que l'ensemble
\[
  A_\lambda = \bigl\{ \{s, \lambda t\} \bigm| \{s, t\} \in A_0 \bigr\}
\]
soit l'ensemble des arêtes d'une structure triangulée \emph{isomorphe} à
$\Sigma_0$. Comme $\lambda$ centralise $G_0$, on a $\mathrm{Aut}(S, A_\lambda) =
G_0$. D'où les définitions (pages 55 à 57) : deux graphes stricts $A_0$, $A_1$
sur $S$ sont \emph{équivalents} s'il existe $\lambda \in \mathfrak{S}_S$ qui
centralise $\mathrm{Aut}(S, A_0)$ et vérifie $\{s, t\} \in A_0 \iff \{s,
\lambda t\} \in A_1$ (on dit que $\lambda$ « transporte » $A_0$ en $A_1$) ; ils
sont \emph{strictement équivalents} s'ils sont de plus isomorphes. La relation
est symétrique ($\lambda^{-1}$ transporte $A_1$ en $A_0$) et transitive
\footnote{Si $\lambda$ transporte $A_0$ en $A_1$ et $\mu$ transporte $A_1$ en
$A_2$, c'est $\mu\lambda$ qui transporte $A_0$ en $A_2$ ; la page écrit
$\lambda\mu$, ce qui revient au même dès que $\lambda$ et $\mu$ commutent,
hypothèse que la suite fait.}. L'ensemble $\widetilde{Z}_0$ des $\lambda$
qui transportent $A_0$ en une structure équivalente est stable par inverse et
par produit d'éléments qui commutent ; c'est un groupe sous l'hypothèse,
faite désormais, que ses éléments commutent deux à deux. Alors les structures
strictement équivalentes à $A_0$ forment un torseur sous $Z'_0 =
Z_0/\mathfrak{z}_0$, $\mathfrak{z}_0 = Z_0 \cap G_0$ (147) ; le sous-groupe $G$
de $\mathrm{Aut}(S, \Pi)$ qui centralise $Z_0$ modulo $\mathfrak{z}_0$
s'insère dans une suite exacte
\[
  1 \longrightarrow G_0 \longrightarrow G \longrightarrow Z'_0 \longrightarrow 1 ;
\]
et si $G_0$ est transitif sur les arêtes, les structures d'une même classe sont
deux à deux disjointes, d'où une application « type » $A = \bigcup A_i \to
\Pi$, qui généralise celle des pages 48 et 49.

Pour l'appliquer au module $M$ sur $\mathbf{Z}/n\mathbf{Z}$, il faudrait
vérifier que $Z_0$ est le groupe des $T_\lambda$, $\lambda \in Z'_n$, et l'on
« suspecte » même que le commutant de $G_0 = \mathrm{GL}^{\pm}(M)'$ tout entier
soit ce groupe (155) ; le § VI le vérifiera pour le commutant de
$\mathrm{SL}(M)'$.

\pagerange{61}{70}
\section*{IV. Anneau quelconque : types, et ce que les arêtes déterminent (pages 61 à 70)}

Soit $\Lambda$ un anneau non nul quelconque. $M^{*} \to S_M$ est un revêtement
de degré $2$ si $2 \neq 0$, de degré $1$ si $2 = 0$. Les arêtes $A_M$ ont un
type $\tau : A_M \to \Pi_M$, compatible à l'action de $G_M = \mathrm{GL}(M)'$,
qui opère sur $\Pi_M$ à travers $\delta_M : G_M \to Z'_\Lambda$, déterminant
modulo $\pm 1$. Le stabilisateur d'un type est
$G^{\circ}_M = \operatorname{Ker} \delta_M$\footnote{La page 63 l'identifie à
$\mathrm{SL}(M)/\{\pm 1\}$. Le noyau est $\mathrm{GL}^{\pm}(M)/\{\pm 1\}$, qui
contient $\mathrm{SL}(M)/\{\pm 1\}$ avec indice $2$ quand $2 \neq 0$.}, transitif
sur chaque $A_{M,i}$ ; $G_M$ est transitif sur $A_M$. Les faces $F_M$ ont un
type, et l'on distingue :
\begin{itemize}
  \item A) $2 = 0$ : toute arête est dans une seule face (car $z = x + y$) ;
    les arêtes orientées forment un torseur sous $G_M$ ;
  \item B) $2 \neq 0$ : toute arête est dans exactement deux faces ; les
    « repères » $(s, t, u)$, $\{s, t, u\} \in F_M$, s'identifient aux bases
    modulo $\pm 1$ par $(x, y) \mapsto (\dot x, \dot y, (x+y)^{\cdot})$ et
    forment un torseur sous $G_M$.
\end{itemize}

Les arêtes déterminent-elles les faces ? Dans le cas B, la page voudrait que
les faces soient exactement les triangles du graphe $(S_M, A_M)$, et écrit
aussitôt « relation fausse !! » : trois éléments deux à deux en base ne
vérifient pas en général $z = \pm x \pm y$\footnote{La page ajoute « même sur
$\mathbf{Z}/n\mathbf{Z}$ ou $\mathbf{Z}$ ». Sur $\mathbf{Z}$ c'est faux : si
$(x, y)$ est une base et $z = ax + by$ avec $x \wedge z$, $y \wedge z$
inversibles, alors $a, b \in \mathbf{Z}^{*} = \{\pm 1\}$. La relation tient de
même pour $\mathbf{Z}/n\mathbf{Z}$, $n \in \{3, 4, 6\}$, où il n'y a qu'un type.
Elle échoue dès que $\Lambda^{*} \neq \{\pm 1\}$ : sur $\mathbf{Z}/5$, $\{x, y,
x + 2y\}$ est un triangle du graphe qui n'est pas une face.}. En revanche, si
l'on connaît aussi les types, les faces de type $i$ sont exactement les
triangles du graphe des arêtes de type $i$ (178) : de $x \wedge y$, $y \wedge
z$, $z \wedge x \in \{\pm e\}$ et $z = \varepsilon x + \varepsilon' y$, on tire
$\varepsilon' x \wedge y = x \wedge z$, d'où $\varepsilon' = \pm 1$, et de même
pour $\varepsilon$.

Dans le cas A (caractéristique $2$), la page croit d'abord que les faces avec
leurs types sont « nettement plus fines » que les arêtes avec leurs
types\footnote{Pages 68 et 69 ; le mot lu « jauge » y est conjectural, et
l'exemple invoqué, « le corps premier $\mathbf{F}_2$ », ne peut pas être celui
visé : sur $\mathbf{F}_2$ toute la structure se réduit à un seul triangle. La
formule (15) de la page 74 réfute d'ailleurs l'idée : quand $2 = 0$, les
triangles du graphe de type $i$ sont exactement les faces de type $i$.}, puis
introduit les faces orientées et leur type à valeurs dans $(\det M)^{*}$ tout
entier (179) à (181), en se demandant dans quelle mesure cette structure se
déduit des autres. La page 70 ne porte que cette question et un titre barré.

\pagerange{71}{80}
\section*{V. Récapitulation (pages 71 à 91)}

\subsection*{La structure $\widetilde{P}_M$ et son groupe (pages 71 à 80)}

Le § V reprend tout sur un anneau commutatif non nul quelconque, avec une
nouvelle numérotation. Sur $\widetilde{S}_M = M^{*}$, les faces
$\widetilde{F}_M$ sont les $\{x, y, z\}$, $x + y + z = 0$, $(x, y)$ base ; une
face orientée $x \to y \to z$ a pour type $x \wedge y \in
\overrightarrow{\Pi}_M = (\det M)^{*}$, le type de la face non orientée est sa
classe dans $\Pi_M$, et les $\widetilde{F}_{M,i}$ forment une famille de
structures triangulées sur le même ensemble $M^{*}$, indexée par $\Pi_M$. Toute
arête (paire non ordonnée formant une base) est contenue dans \emph{une seule}
face : $\widetilde{P}_M$ n'est pas un 2-polyèdre, mais son quotient par $\pm 1$
en sera un. Il faudra interpréter $\Pi_M$ comme torseur sous $Z'_\Lambda$, et
$\overrightarrow{\Pi}_M$ comme torseur sous $\Lambda^{*}$ qui le relève (« tirer
au clair »). Les triangles du graphe de type $i$ sont les $\{x, y, \pm x \pm
y\}$, de sorte que les arêtes typées ne déterminent les faces que si $2 = 0$
(15).

L'\emph{antipodie} est l'automorphisme induit par $-\mathrm{id}_M$. Un
automorphisme $u$ de $\widetilde{P}_{M,i}$ qui commute à l'antipodie et fixe
$x$ et $y$, avec $x \wedge y \in i$, fixe aussi $-(x+y)$ ; de proche en proche,
l'ensemble des repères fixés par $u$ est stable par composition avec les
matrices $\tau_0^r$, $\tau_1^r$, $\tau_\infty^r$ du § II, qui engendrent
l'image de $\mathrm{SL}(2, \mathbf{Z})$, c'est-à-dire $\mathrm{SL}(2,
\Lambda_0)$. D'où :

\emph{Lemme} (page 76). Un tel $u$ fixe tous les $ax + by$ avec $a, b \in
\Lambda_0$ et $(a, b)$ unimodulaire dans $\Lambda_0$\footnote{La page écrit
« premiers entre eux (i.e. $a\Lambda + b\Lambda = \Lambda$, i.e. $\exists c, d
\in \Lambda$ avec $ad - bc = 1$) ». L'argument donne des $c, d$ dans
$\Lambda_0$, et l'énoncé avec $c, d \in \Lambda$ est trop fort : sur $\Lambda =
\mathbf{Q}$, $\Lambda_0 = \mathbf{Z}$, le couple $(2, 0)$ est unimodulaire
dans $\mathbf{Q}$, et rien dans la démonstration ne dit que $u$ fixe $2x$.}.

\emph{Corollaire.} Si $\Lambda = \Lambda_0$, un automorphisme de
$\widetilde{P}_{M,i}$ qui fixe un repère et commute à l'antipodie est
l'identité.

\emph{Théorème} (page 77). Soit $\Lambda \simeq \mathbf{Z}$ ou
$\mathbf{Z}/n\mathbf{Z}$. a) Pour $i \in \Pi_M$, $\mathrm{GL}^{\pm}(M) \to
\mathrm{Aut}(\widetilde{P}_{M,i})^{\mathrm{antipod}}$ (automorphismes commutant
à l'antipodie) est un isomorphisme, car $\mathrm{GL}^{\pm}(M)$ est simplement
transitif sur les bases de type $i$. b) $\mathrm{GL}(M) \simeq
\mathrm{Aut}(\widetilde{P}_M)$, la structure étant la famille
$(\widetilde{F}_{M,i})_i$ avec l'antipodie. Ici il n'y a pas de quotient par
$\pm 1$ : l'antipodie fait partie de la structure.

\emph{Scholie} (page 78). Par les principes généraux, on a deux équivalences
de groupoïdes, pour $\Lambda_n = \mathbf{Z}$ ($n = 0$) ou
$\mathbf{Z}/n\mathbf{Z}$ : $(M, i) \mapsto \widetilde{P}_{M,i}$, des modules
libres de rang $2$ munis d'une structure semi-symplectique vers les structures
triangulées munies d'une involution (l'antipodie) isomorphes à une structure
type dépendant de $n$ ; et $M \mapsto \widetilde{P}_M$, des modules libres
vers les structures $(S, F, \Pi, \tau, a)$. « Il conviendrait de caractériser
axiomatiquement, par des propriétés intrinsèques », les structures du second
membre, et de comprendre comment reconstruire $M$.

\pagerange{79}{83}
\subsection*{L'action de $\Lambda^{*}$, et la présentation de $M$ (pages 79 à 83)}

Quand $\Lambda \neq \Lambda_0$, on ajoute l'action $\widetilde{T}_\lambda x =
\lambda x$ de $\Lambda^{*}$ sur $M^{*}$ : elle respecte $\widetilde{F}_M$ et la
partition en types, et commute à $\mathrm{GL}(M)$, dont elle est le centre. Un
automorphisme $u$ de la structure qui commute aux $\widetilde{T}_\lambda$ et
fixe une base $(x, y)$ fixe les $\lambda y$, donc tous les types $\lambda i$, et
respecte chaque $\widetilde{F}_{M,i}$. Si $\Lambda^{*}$ engendre $\Lambda$
additivement (par exemple $\Lambda$ semi-local), les matrices
$\bigl(\begin{smallmatrix} 1 & -\lambda \\ 0 & 1 \end{smallmatrix}\bigr)$,
$\lambda \in \Lambda^{*}$, avec $\lambda_1$ et $\lambda_\infty$, engendrent
« sauf erreur » $\mathrm{SL}(2, \Lambda)$, et l'on trouve que tout automorphisme
de $(\widetilde{S}_M, \widetilde{F}_M, \widetilde{T})$ provient d'un
automorphisme unique de $M$\footnote{Le « sauf erreur » se lève : pour un
anneau semi-local, $\mathrm{SL}(2, \Lambda)$ est engendré par les matrices
élémentaires (résultat classique, cité de mémoire), et une matrice élémentaire
de coefficient $t = \sum_k \lambda_k$, $\lambda_k \in \Lambda^{*}$, est le
produit de celles de coefficients $\lambda_k$.}.

En termes additifs, une permutation $u$ de $M^{*}$ qui commute à l'antipodie et
respecte les faces vérifie $u(-x) = -u(x)$ et $u(x + y) = u(x) + u(y)$ dès que
$(x, y)$ est une base ; commuter à $\Lambda^{*}$ ajoute $u(\lambda x) = \lambda
u(x)$. Il est « tentant » d'écrire $M$ comme le quotient du groupe abélien libre
$\mathbf{Z}^{(M^{*})}$ par le sous-groupe $R_M$ engendré par les $[x + y] - [x] -
[y]$ ($(x, y)$ base) et les $[-x] + [x]$ : $u$ s'étendrait alors en un
automorphisme additif $\hat u$ de $M$. Il faudrait montrer que la suite $0 \to
R_M \to \mathbf{Z}^{(M^{*})} \to M \to 0$ est exacte\footnote{Pour $\Lambda =
\Lambda_0$ c'est vrai, et l'argument est celui du lemme de la page 76 (il est
de nous) : dans le quotient, $[x]$ et $[y]$ commutent ($[x] + [y] = [x+y] =
[y] + [x]$), $[kx + y] = k[x] + [y]$ par récurrence, d'où $n[x] = 0$ sur
$\mathbf{Z}/n\mathbf{Z}$ ; et tout élément de $M^{*}$, de la forme $ax + by$ avec
$(a, b)$ unimodulaire, s'obtient à partir de $\pm x$, $\pm y$ par des additions
de bases (pas d'Euclide). Le quotient est donc engendré par $[x]$, $[y]$ et se
surjecte sur $M$ : il lui est isomorphe.} ; une note marginale y voit le groupe
« enveloppant » d'une loi de groupe partiellement définie.

Quand $\Lambda$ est « absolu », tout endomorphisme additif est
$\Lambda$-linéaire, et la condition sur $\Lambda^{*}$ est superflue. La page
explique pourquoi elle est indispensable en général : « tout automorphisme
\emph{additif} (non nécessairement $\Lambda$-linéaire) de $M$ induit un
automorphisme de $\widetilde{P}_M$ »\footnote{Pas tout automorphisme additif :
sur $\Lambda = \mathbf{F}_4 = \mathbf{F}_2(\omega)$, l'automorphisme
$\mathbf{F}_2$-linéaire de $\mathbf{F}_4^2$ qui fixe $(1, 0)$ et $(0, \omega)$,
envoie $(\omega, 0)$ sur $(0, 1)$ et $(0, 1)$ sur $(\omega, 1)$, envoie la base
$((1, 0), (\omega, 1))$ sur la paire liée $((1, 0), (\omega, 0))$. Ce qui est
vrai, et suffit à l'argument, c'est que tout automorphisme \emph{semi-linéaire}
(linéaire au-dessus d'un automorphisme de l'anneau) respecte la structure.}.
La partition en types l'empêche dans certains cas, mais non si $\Lambda^{*} =
\{\pm 1\}$, par exemple pour un anneau de polynômes sur $\mathbf{F}_2$,
$\mathbf{F}_3$ ou $\mathbf{Z}$, où il n'y a qu'un type.

\pagerange{84}{91}
\subsection*{Le quotient par $\pm 1$ (pages 84 à 91)}

On suppose $2 \neq 0$. Alors $-\mathrm{id}_M$ opère librement sur $M^{*}$, sur
les arêtes et sur les faces ; $S_M$, $F_M \subset \mathfrak{P}_3(S_M)$ et les
types se définissent par passage au quotient. Au-dessus d'une arête $\{\dot x,
\dot y\}$ il y a quatre arêtes de $\widetilde{P}_M$, $\{x, y\}$, $\{-x, -y\}$,
$\{x, -y\}$, $\{-x, y\}$, groupées en deux paires antipodiques qui
correspondent aux deux faces contenant l'arête (35). Une jauge $e$ choisie dans
$i = \{\pm e\}$ oriente $P_{M,i}$, de sorte que l'ensemble des orientations de
$P_{M,i}$ s'identifie à $i$ ; les faces de type $i$ se reconstituent à partir des
arêtes de type $i$ (37). L'étoile de $\dot x$ dans $P_{M,i}$ est la droite
affine $L_x = \{y \mid x \wedge y = e\}$, où le successeur de $y$ est $y + x$
(39), d'où (40) : les étoiles sont connexes si et seulement si $\Lambda =
\Lambda_0$\footnote{La page ajoute que, dans tous les cas, les « fragments »
(composantes) de l'étoile sont d'ordre $\operatorname{card}\Lambda$. Ce sont
les classes $y + \Lambda_0 x$, de cardinal $\operatorname{card}\Lambda_0$ ; les
deux nombres ne coïncident que si $\Lambda = \Lambda_0$.}, et $P_M$ tout entier
est un 2-polyèdre si et seulement si de plus il n'y a qu'un type, $\Lambda
\simeq \mathbf{Z}$ ou $\mathbf{Z}/n\mathbf{Z}$, $n \in \{3, 4, 6\}$ (41).

$\mathrm{GL}^{\pm}(M)'$ est simplement transitif sur les repères de $P_{M,i}$,
$\mathrm{SL}(M)'$ sur les repères directs de $P_{M,e}$, $\mathrm{GL}(M)'$ sur
les repères de $P_M$\footnote{La page 89 écrit le dernier groupe
« $\mathrm{Gl}'_\Lambda(M)^{\varepsilon}$ » ; il ne peut s'agir que de
$\mathrm{GL}(M)'$.}. D'où :

\emph{Théorème} (page 89). Si $\Lambda = \Lambda_0$ : a) $\mathrm{SL}(M)'
\simeq \mathrm{Aut}^{+}(P_{M,e})$ ; b) $\mathrm{GL}^{\pm}(M)' \simeq
\mathrm{Aut}(P_{M,i})$ ; c) $\mathrm{GL}(M)' \simeq \mathrm{Aut}(P_M)$, où
$\mathrm{Aut}(P_M)$ est le groupe des permutations de $S_M$ qui respectent
l'ensemble $F_M$ de toutes les faces\footnote{Pour c), une démonstration
(de nous) : un automorphisme de $P_M$ permute les $P_{M,i}$ (une note marginale
remarque que ce sont les composantes connexes du graphe d'adjacence des
faces) ; on le compose avec un élément de $\mathrm{GL}(M)'$ qui ramène l'image
de $P_{M,i}$ sur $P_{M,i}$, puis on applique b). L'image de $\mathrm{Aut}(P_M)$
dans les permutations des types n'est que $Z'_\Lambda$, opérant par
translations ; la légende de la suite exacte (42) le dit (« pas
surjectif »).}.

C'est le théorème de la page 42, débarrassé de ses structures supplémentaires :
les polyèdres $P_{M,i}$ étant posés sur le \emph{même} ensemble $S_M$, les
identifications (101) du § III deviennent des identités, et il n'y a plus rien
à comprendre « a priori ».

\emph{Scholie} (pages 89 à 91). Pour $\Lambda \simeq \mathbf{Z}$ ou
$\mathbf{Z}/n\mathbf{Z}$ fixé, on a trois équivalences de groupoïdes : modules
symplectiques libres de rang $2$ modulo $\pm 1$ vers 2-polyèdres triangulés
orientés réguliers d'un type déterminé par $n$ ; modules semi-symplectiques vers
2-polyèdres réguliers (orientables, non orientés) ; modules libres vers
« structures multi-2-polyédrales triangulées ». Le programme qui suit occupera
le reste du dossier : a) caractériser axiomatiquement ces structures ; b) les
expliciter ; c) reconstruire $M$ ; d) expliciter les structures annexes (le
torseur $\Pi$, le $\Lambda^{*}$-torseur $\overrightarrow{\Pi}$ des orientations,
l'action de $Z'_\Lambda$ sur $S_M$) ; e) expliquer comment toutes les $P_{M,i}$
se déduisent de l'une ; f) pour $\Lambda$ quelconque, trouver les structures
combinatoires qui portent la trace de $\Lambda$.

\pagerange{92}{95}
\section*{VI. Le commutant, et $Z'_\Lambda$ lu sur les étoiles (pages 92 à 101)}

\subsection*{Le commutant de $\mathrm{SL}(M)$ et de $\mathrm{SL}(M)'$ (pages 92 à 95)}

La page commence par la question e). Si $S$ est un espace homogène sous un
groupe $G$ et $H$ le stabilisateur d'un point, le commutant de $G$ dans
$\mathfrak{S}_S$ est isomorphe à $N_G(H)/H$\footnote{Fait classique ; la page
le rappelle sans démonstration. Il revient au § X, page 179, avec la
convention de la translation à droite.}. On calcule ce normalisateur dans une
base $(x, y)$.

A) $G = \mathrm{SL}(2, \Lambda)$ opérant sur $M^{*}$, $H = \bigl\{
\bigl(\begin{smallmatrix} 1 & t \\ 0 & 1 \end{smallmatrix}\bigr) \bigr\}$ le
stabilisateur de $x$. Pour $g = \bigl(\begin{smallmatrix} a & b \\ c & d
\end{smallmatrix}\bigr)$,
\[
  g \begin{pmatrix} 1 & t \\ 0 & 1 \end{pmatrix} g^{-1} =
  \begin{pmatrix} 1 - act & a^2 t \\ -c^2 t & 1 + act \end{pmatrix} ,
\]
qui est dans $H$ pour tout $t$ si et seulement si $c^2 = 0$ et $ac = 0$ ; alors
$ad = 1 + bc$ est inversible, donc $a$ aussi, et $c = 0$. Le normalisateur est
le groupe des matrices triangulaires supérieures $\bigl(\begin{smallmatrix}
\lambda & b \\ 0 & \lambda^{-1} \end{smallmatrix}\bigr)$\footnote{La page 93
écrit en passant « $(1 - ac\lambda)(1 + ac\lambda) = 1$ », sans suite. Le
groupe trouvé est ce qu'on appelle aujourd'hui un sous-groupe de Borel ; le nom
n'est pas sur la page, mais les gloses « unipotents » et « Borel » reviennent
plus loin (page 153).}, et le commutant de $\mathrm{SL}(M)$ dans les
permutations de $M^{*}$ est formé des homothéties $\lambda\,\mathrm{id}_M$,
$\lambda \in \Lambda^{*}$.

B) $G' = \mathrm{SL}(2, \Lambda)/\{\pm 1\}$ opérant sur $S_M$, $2 \neq 0$. La
condition devient : $c^2 t = 0$ et $1 - act = \pm 1$ pour tout $t$. Si $ac = 0$,
on conclut comme en A. Sinon $ac = 2$, et l'on trouve $4 = 0$, $2t \in \{0,
2\}$ pour tout $t$ --- l'idéal $2\Lambda$ n'a que deux éléments --- et $c = 2$.
D'où deux cas :
\begin{itemize}
  \item 1) si $2\Lambda = \{0, 2\} \neq \{0\}$, le normalisateur est formé des
    classes des matrices avec $c \in \{0, 2\}$, et contient le groupe
    triangulaire avec l'indice $2$ ; le commutant de $\mathrm{SL}(M)'$ dans
    $\mathfrak{S}_{S_M}$ est le produit direct de $Z'_\Lambda$ et d'un groupe
    $\{1, \sigma\}$ d'ordre $2$, $\sigma$ étant induit par
    $\bigl(\begin{smallmatrix} 1 & 0 \\ 2 & 1 \end{smallmatrix}\bigr)$ ;
  \item 2) dans tous les autres cas, le commutant de $\mathrm{SL}(M)'$ est
    $Z'_\Lambda$.
\end{itemize}
Le cas typique de 1) est $\Lambda = \mathbf{Z}/4\mathbf{Z}$ : $Z'_\Lambda$ est
trivial, mais les permutations des six sommets de l'octaèdre qui commutent aux
rotations ne se réduisent pas à l'identité --- il y a l'antipodie de
l'octaèdre.

\pagerange{96}{101}
\subsection*{Lire $Z'_\Lambda$ sur les étoiles (pages 96 à 101)}

Soit $\Lambda = \mathbf{Z}/n\mathbf{Z}$, $n \neq 2, 4$, et $P_{i_0}$ l'une des
structures $P_{M,i}$ sur $S = S_M$. Le groupe $G = \mathrm{Aut}(P_{i_0})$ et
son sous-groupe direct $G^{+}$ ne dépendent pas de $i_0$ ; soit $Z$ le
commutant de $G^{+}$ dans $\mathfrak{S}_S$ --- c'est $Z'_n$ par ce qui précède.
Pour $\lambda \in Z$, on a $\{s, t\} \in A_{i_0} \iff \{\lambda s,
\lambda^{-1} t\} \in A_{i_0}$, et l'unique structure dont les arêtes sont les
$\{s', t'\}$ avec $\{s', \lambda^{-1} t'\} \in A_{i_0}$ est l'une des
$A_i$, notée $A_{\lambda \cdot i_0}$ ; l'application $\lambda \mapsto \lambda
\cdot i_0$ est bijective, et l'ensemble des structures est un torseur sous
$Z$. C'est la torsion d'une seule extrémité, cette fois avec le bon groupe.

Il reste à expliciter $Z \simeq Z'_n$ en termes de la seule combinatoire. Pour
$s \in S$ et $\lambda \in Z$, la permutation $t \mapsto \lambda t$ envoie
l'étoile de $s$ dans $P_{i_0}$ sur l'étoile de $s$ dans $P_{i_1}$, $i_1 =
\lambda \cdot i_0$ (61). Chacune de ces étoiles est un torseur sous le groupe
$\Gamma_i = \Lambda \wedge_{\mu_2} \mu(i)$, le groupe additif $\Lambda$ tordu
par l'ensemble $\mu(i)$ des deux orientations de $P_i$ : choisir une
orientation, c'est choisir le sens de la translation $y \mapsto y + x$. La page
affirme (« à vérifier ») que $\lambda_s$ est un isomorphisme de torseurs pour
un isomorphisme de groupes $\lambda_{\Gamma, s} : \Gamma_{i_0} \to
\Gamma_{i_1}$, c'est-à-dire, une fois les orientations choisies, la
multiplication par un élément de $\Lambda^{*}$ défini au signe près, et que
l'application $Z \to Z'_\Lambda$ ainsi obtenue ne dépend ni de $s$ ni de $i_0$
et est l'isomorphisme inverse de celui de la page 47\footnote{Le calcul (de
nous) confirme l'énoncé. Pour $\lambda = T_\mu$, multiplication par $\mu \in
\Lambda^{*}$, et $s = \dot x$ : l'étoile de $s$ dans $P_{i_0}$ est $\{y + kx\}$,
avec les voisins $y \sim y \pm x$ ; son image est $\{\mu y + k\mu x\}$, étoile
de $s$ dans $P_{\mu i_0}$, où les voisins sont $y' \sim y' \pm x$. La
translation de $k$ crans devient donc une translation de $\mu k$ crans :
$\lambda_{\Gamma, s} = \pm\mu$, et l'application est $T_\mu \mapsto \mu$ modulo
$\pm 1$, sans exposant. Les formules (65) et (66) portent un exposant illisible
(« un $4$ ? ») qu'on ne restitue pas.}.

Enfin, l'ensemble $\overrightarrow{\Pi} = \coprod_{i \in \Pi} \mu(i)$ des
couples (structure, orientation) reçoit une action de $\Lambda^{*}$ au-dessus
de celle de $Z'_\Lambda$ sur $\Pi$, construite de même à partir de
$\lambda_{\Gamma, s}$ : il existe une unique orientation $e'$ de $P_{\lambda i}$
telle que $\lambda_{\Gamma, s} = \lambda \wedge (e \wedge e')$, et ce choix ne
dépend pas de $s$. Cela fait de $\overrightarrow{\Pi}$ un torseur sous
$\Lambda^{*}$\footnote{La page 101 écrit « torseur sous $\Lambda$ » ; il s'agit
de $\Lambda^{*}$, qui relève $Z'_\Lambda = \Lambda^{*}/\{\pm 1\}$.}, et le type
orienté des faces orientées $\overrightarrow{F} \to \overrightarrow{\Pi}$
relève le type non orienté.

\pagerange{102}{106}
\section*{VII. Structures triangulaires congruentielles (pages 102 à 128)}

\subsection*{Structures triangulaires et leurs composantes (pages 102 à 106)}

Le titre est celui de la page 102 ; le mot « congruentielles » n'y est pas
expliqué. Une \emph{structure triangulaire} est un couple $(S, F)$, $F \subset
\mathfrak{P}_3(S)$, avec (T1) $S = \bigcup_{f \in F} f$. Avec $A$ l'ensemble des
paires contenues dans une face, on demande (T2) : toute arête est dans
exactement deux faces. Le graphe $\Gamma$ dont les sommets sont les faces et
les arêtes les paires de faces adjacentes est alors un graphe strict à sommets
d'ordre $3$ ; l'ensemble $\pi = \pi_0(\Gamma)$ de ses composantes partitionne
$F$ et $A$ en $F_i$, $A_i$ ($i \in \pi$). Les relations d'incidence vérifient
l'axiome des structures 2-polyédrales (toute arête a deux sommets et deux
faces ; un sommet et une face incidents ont exactement deux arêtes
communes), mais l'étoile d'un sommet $P(s) = \coprod_i P_i(s)$ n'est pas
connexe. On ne le demande pas ; on demande plutôt (T3) : pour tout $i$, $(S,
F_i)$ est une structure triangulaire \emph{stricte} --- ses faces recouvrent
$S$, et ses étoiles sont connexes (T3 a, b, page 106). Une telle structure
s'interprète comme un ensemble $\pi$ de 2-polyèdres triangulés connexes sur le
même ensemble de sommets ; c'est la forme abstraite de $P_M$.

\pagerange{107}{112}
\subsection*{Les axiomes indexés par $0$ (pages 107 à 112)}

On part maintenant d'\emph{un} 2-polyèdre $P_0 = (S, A_0, F_0)$ et l'on cherche
des axiomes qui permettent de reconstruire tous les autres :
\begin{itemize}
  \item ($\mathrm{T}_0 1$) $P_0$ est \emph{régulier} : $G_0 =
    \mathrm{Aut}(P_0)$ est transitif sur les repères $(s, t, u)$, $\{s, t,
    u\} \in F_0$ ;
  \item ($\mathrm{T}_0 2$) $P_0$ est orientable ; $G_0^{+}$, d'indice $2$, est
    le groupe des automorphismes directs, simplement transitif sur les arêtes
    orientées\footnote{La page 108 écrit « $P_0$ est simplement transitif » ;
    il s'agit de $G_0^{+}$.}.
\end{itemize}
Soit $Z' \subset \mathfrak{S}_S$ le commutant de $G_0^{+}$ ; $S$ étant un espace
homogène, $Z' \simeq N/G^{+}_{0s}$, $N$ le normalisateur dans $G_0^{+}$ du
stabilisateur $G^{+}_{0s}$ d'un sommet. Ce stabilisateur est cyclique d'ordre
$n$ (l'ordre des sommets) ; canoniquement, $G^{+}_{0s} \simeq \Lambda_n
\wedge_{\mu_2} \mu_0$, le groupe $\Lambda_n = \mathbf{Z}/n\mathbf{Z}$ tordu par
l'ensemble $\mu_0$ des deux orientations\footnote{Marge : « le sens des deux
symboles $\mu_2$ et $\mu_0$ est tout à fait différent ! » : $\mu_2 = \{\pm 1\}$
est le groupe, $\mu_0$ un ensemble à deux éléments.}. La conjugaison donne
donc un caractère $\chi_s : Z' \to \Lambda_n^{*}$. Mais dans le modèle, $\chi_s(\dot\lambda)
= \lambda^2$ : il n'est pas injectif, « il faut trouver un autre principe pour
construire un iso $Z' \simeq \Lambda_n^{*}/\pm 1$ »\footnote{Dans le modèle, le
normalisateur du groupe unipotent $\bigl(\begin{smallmatrix} 1 & t \\ 0 & 1
\end{smallmatrix}\bigr)$ contient $\mathrm{diag}(\lambda, \lambda^{-1})$, qui
agit par $t \mapsto \lambda^2 t$ : c'est le $\lambda^2$ des conventions, et
c'est l'obstacle sur lequel le dossier reviendra jusqu'à ses dernières pages
(pages 292 à 294).}.

Pour $\lambda \in Z'$ et $s \in S$, la paire $\{s, \lambda s\}$ n'est jamais une
arête (31) : sinon, soit $t$ le troisième sommet d'une face sur cette arête, et
$\rho \in G_0^{+}$ la rotation autour de $\lambda s$ qui envoie $s$ sur $t$ ;
de $\rho(\lambda s) = \lambda \rho(s) = \lambda s$ on tirerait $\rho(s) = s$,
donc $t = s$. On suppose alors
\begin{itemize}
  \item ($\mathrm{T}_0 3$) pour $\lambda \in Z'$ : $\{s, t\} \in A_0 \iff
    \{\lambda s, \lambda^{-1} t\} \in A_0$.
\end{itemize}
Il suffit de le vérifier pour un repère origine $(s_0, t_0, u_0)$, et l'axiome
se traduit en théorie des groupes ($\mathrm{T}_0 3'$, page 111) : avec $t_0 =
\rho s_0$, pour tout $g$ normalisant $G^{+}_{0 s_0}$, il existe $z \in G^{+}_{0
s_0}$ (nécessairement unique) tel que $g \cdot \mathrm{int}(\rho^{-1})(gz) \in
G^{+}_{0s_0}$. Un 2-polyèdre orienté régulier épinglé équivaut en effet à la
donnée du groupe $G_0^{+}$ avec ses générateurs $\rho_0, \rho_1, \rho_\infty$,
$\rho_0^n = \rho_1^2 = \rho_\infty^3 = 1$, $G^{+}_{0 s_0}$ étant engendré par
$\rho_0$\footnote{Ces relations ne sont pas une présentation : il faut ajouter
$\rho_\infty\rho_1\rho_0 = 1$ (page 27), et $G_0^{+}$ est un \emph{quotient} du
groupe triangulaire ainsi présenté. La page 143 le redira.}. Si l'on dispose de
$G_0$ tout entier, on peut prendre pour $\rho$ la réflexion $\tau_0$ qui échange
$s_0$ et $t_0$ (34').

\pagerange{113}{120}
\subsection*{Le graphe tordu $\Gamma_\lambda$ (pages 113 à 120)}

Pour $\lambda \in Z'$, l'ensemble $A_\lambda = \{\{s, \lambda t\} \mid \{s, t\}
\in A_0\}$ est un graphe strict $\Gamma_\lambda$ (par (31), $s \neq \lambda
t$) ; $\lambda$ induit une bijection de l'étoile de $s$ dans $\Gamma_0$ sur
son étoile dans $\Gamma_\lambda$, de sorte que tous les sommets de
$\Gamma_\lambda$ sont d'ordre $n$, et $G_0^{+}$ opère sur $\Gamma_\lambda$,
simplement transitivement sur les arêtes orientées. Pour que $\Gamma_\lambda$
soit le graphe d'un 2-polyèdre régulier, la page ajoute :
\begin{itemize}
  \item ($\mathrm{T}_0 4$) le graphe détermine les faces : $F_0$ est l'ensemble
    des triangles de $A_0$ ;
\end{itemize}
axiome « par acquit de conscience », dont la marge demande s'il suit des
précédents. « J'ignore s'il existe une structure de 2-polyèdre régulier
triangulé orientable qui ne soit [pas] décrite par sa structure de graphe »,
et, sans l'hypothèse d'orientabilité, il en donne un : l'icosaèdre « gauche »,
régulier non orientable, à $6$ sommets, dont les $15$ paires sont toutes des
arêtes, mais dont une partie seulement des $20$ triplets sont des
faces\footnote{C'est l'\emph{hémi-icosaèdre}, quotient de l'icosaèdre par
l'antipodie, tracé sur le plan projectif ; il a $10$ faces (le nombre est
illisible sur la page ; il est de nous). Voir le dossier 86, page 28. Pour le
cas orientable, la réponse est aussi négative (de nous) : la carte régulière
$\{3, 6\}_{(3,0)}$ du tore --- les sommets $(\mathbf{Z}/3)^2$, voisins de $0$
les $\pm(1,0)$, $\pm(0,1)$, $\pm(1,1)$ --- a $9$ sommets et $18$ faces, et le
triangle $\{0, (1,0), (2,0)\}$ de son graphe n'est pas une face. Dans le modèle,
en revanche, $\mathrm{T}_0 4$ est vrai : c'est la formule (178) de la page 67.}.
Un graphe définit un 2-polyèdre si (A1) toute arête $\{s, t\}$ a exactement
deux sommets voisins de ses deux extrémités et (A2) l'étoile de chaque sommet,
avec les arêtes qu'elle contient, est un contour connexe. Pour $\Gamma_\lambda$,
(A1) s'écrit
\[
  \operatorname{card}\bigl(\lambda \cdot \mathrm{Et}_0(s) \cap \mathrm{Et}_0(t)\bigr)
  = 2 \qquad \text{si } \{s, t\} \in A_0 ,
\]
ce qui « n'a pas l'air tautologique ». La page pose donc l'axiome substantiel :
\begin{itemize}
  \item ($\mathrm{T}_0 5$) pour tout $\lambda \in Z'$, $\Gamma_\lambda$ est
    isomorphe à $\Gamma_0$ : il existe $g \in \mathfrak{S}_S$ avec $\{s, t\}
    \in A_0 \iff \{gs, \lambda^{-1} gt\} \in A_0$ ;
\end{itemize}
d'où un 2-polyèdre régulier $P_\lambda \simeq P_0$, et
\begin{itemize}
  \item ($\mathrm{T}_0 6$) $Z'$ commute à $G_0$ tout entier,
\end{itemize}
de sorte que $G_0$ opère sur $P_\lambda$ ; l'injection $G_0 \to
\mathrm{Aut}(P_\lambda)$ est bijective quand $S$ est fini, par cardinalité. Que
$G_0^{+}$ soit aussi le groupe direct de $P_\lambda$ est clair si $n$ est
impair ($G_0^{+}$ est le seul sous-groupe d'indice $2$) ; si $n$ est pair, il
y a a priori trois sous-groupes d'indice $2$, l'abélianisé du groupe
cartographique étant $\mathbf{F}_2 \times \mathbf{F}_2$, mais un seul contient
$\rho_1$ --- « il faudra y revenir ailleurs »\footnote{La question est réglée
aux pages 149 et 150 : c'est le seul qui contient la rotation $\rho_0$ autour
d'un sommet, et $\rho_0$ reste une rotation pour $P_\lambda$.}. Les structures
qui satisfont $\mathrm{T}_0 1$ à $\mathrm{T}_0 6$ sont dites « spéciales ».

\pagerange{120}{125}
\subsection*{Équivalence, et l'exception octaédrique (pages 120 à 125)}

Deux structures spéciales sur $S$ sont \emph{équivalentes} s'il existe $\lambda$
commutant à leur groupe tel que $\{s, t\} \in A_0 \iff \{s, \lambda t\} \in
A_1$ ; c'est une relation d'équivalence, et à une classe $\pi$ correspondent un
groupe $G_0$ commun et le groupe $Z' = \mathrm{Centr}_{\mathfrak{S}_S}(G_0)$,
transitif sur $\pi$. La page note la relation « plausible » : $Z'$ est
commutatif. Le calcul donne seulement que $(\lambda\mu)s$ et $(\mu\lambda)s$
ont la même étoile ; d'où la commutativité dès qu'un sommet est déterminé par
son étoile. L'exception connue est l'octaèdre, où deux sommets opposés ont la
même étoile ; alors $Z' = \{1, \varepsilon\}$, $\varepsilon$ l'antipodie,
unique élément central non trivial de $G_0 \simeq \mathfrak{S}_4 \times \{\pm
1\}$. Pour les mêmes raisons, $\pi$ est un torseur sous $Z'$ (le stabilisateur
d'une structure est trivial), sauf dans ce cas.

Ces pages, très surchargées, envisagent aussi la « double pyramide » (le double
cône sur un $n$-gone, où les deux pointes ont la même étoile), avec $Z' \simeq
\{1, \sigma\} \times \{1, \varepsilon\}$ pour $n$ pair $\geq 6$, et, pour $n
\equiv 0 \pmod 4$, $n \geq 8$, deux structures déduites l'une de l'autre par
« multiplication par $\frac{n}{2} - 1$ » sur l'équateur\footnote{La lecture est
mince et on ne restitue rien de plus. Deux remarques de nous : la double
pyramide n'est régulière que pour $n = 4$ (c'est l'octaèdre), elle relève donc
des axiomes affaiblis de la page 123 (« $\mathrm{T}_0 3$, $\mathrm{T}_0 4$ »)
et non de $\mathrm{T}_0 1$ ; et $\frac{n}{2} - 1$ est un inversible d'ordre $2$
modulo $n$ exactement quand $4 \mid n$, ce qui explique la condition
$n \equiv 0 \pmod 4$.}.

\pagerange{125}{128}
\subsection*{Le groupe $G$ et l'axiome To 7 (pages 125 à 128)}

Soit $G$ le groupe des permutations de $S$ qui préservent la classe $\pi$. Le
stabilisateur $G_0$ d'un élément de $\pi$ est le noyau de $G \to
\mathfrak{S}_\pi$, donc distingué ; $G$ normalise le centralisateur $Z'$ de
$G_0$, et l'on a $G/G_0 \to \mathrm{Aff}_{Z'}(\pi)$. De plus $Z' \subset G$,
car
\[
  \lambda(A_{i_0}) = A_{\lambda^2 \cdot i_0}
\]
(une marge note : « $\det(\lambda\,\mathrm{id}_M) = \lambda^2$, $\neq \lambda$
en général ! »). On pose
\begin{itemize}
  \item (To 7) $Z'$ est central dans $G$ (donc égal à son centre),
\end{itemize}
ce qui équivaut à dire que $G$ opère sur $\pi$ par translations, via un
homomorphisme $\delta : G \to Z'$ dont la restriction à $Z'$ est $\lambda
\mapsto \lambda^2$. L'image de $\delta$ contient les carrés ; elle est tout
$Z'$ si et seulement si toutes les structures équivalentes sont isomorphes,
c'est-à-dire sous $\mathrm{T}_0 5$, et l'on a alors une suite exacte
\[
  1 \longrightarrow G_0 \longrightarrow G \xrightarrow{\ \delta\ } Z'
  \longrightarrow 1 .
\]
Dans le modèle, $G = \mathrm{GL}(M)'$, $G_0 = \mathrm{GL}^{\pm}(M)'$, $Z' =
Z'_n$ et $\delta$ est le déterminant modulo $\pm 1$, qui vaut $\lambda^2$ sur
les scalaires.

\pagerange{129}{136}
\section*{VIII. Axiomatique : reformulation (pages 129 à 158)}

\subsection*{Structures prégroupales (pages 129 à 136)}

Nouvelle mouture, nouveaux axiomes. Une structure triangulaire $F_0$ sur $S$
doit satisfaire (To 1) : $F_0$ est l'ensemble des faces d'un 2-polyèdre
combinatoire connexe --- toute arête dans exactement deux faces, connexité
globale et locale --- et (To 2) : $F_0$ est l'ensemble des triangles de son
graphe $A_0$. La régularité n'est plus demandée. Avec $G_0 = \mathrm{Aut}(S,
F_0)$, $\widehat{Z}$ son centralisateur, on considère l'ensemble $Z^{*}$ des
$\lambda \in \widehat{Z}$ tels que a) $\{s, t\} \in A_0 \iff \{\lambda^{-1} s,
\lambda t\} \in A_0$ --- ce qui définit $A_\lambda$ --- et b) $\Gamma_\lambda =
(S, A_\lambda)$ est isomorphe à $\Gamma_0$. On demande (To 3) que les éléments
de $Z^{*}$ commutent deux à deux, ce qui équivaut à ce que $Z^{*}$ soit un
sous-groupe $Z$, nécessairement commutatif. Alors $\pi = \{A_\lambda\}$ est
un ensemble homogène sous $Z$, $G = \{f \mid f(A_0) \in \pi\}$, $\mathrm{Cent}(G)
\subset Z \subset G \supset G_0$, $G/G_0 \hookrightarrow \mathrm{Aff}_Z(\pi)$,
et l'on demande (To 4) que $Z$ soit central dans $G$, d'où la suite exacte
canonique
\[
  1 \longrightarrow G_0 \longrightarrow G \xrightarrow{\ \chi'\ } Z
  \longrightarrow 1 ,
\]
$\chi'$ valant $\lambda \mapsto \lambda^2$ sur $Z$ (dans le cas de la double
pyramide, il faut remplacer $Z$ par $Z/\{1, \varepsilon\}$). Ces structures
sont dites « prégroupales symplectiques » ; deux d'entre elles sont
équivalentes (18) si elles sont isomorphes et si l'une se déduit de l'autre par
torsion d'une extrémité au moyen d'une permutation commutant à $G_0$ ; une
« structure triangulaire prégroupale » est une classe d'équivalence. La page
reconnaît que l'axiomatique n'est « pas encore réduite au minimum », et
qu'elle se dissiperait peut-être plus naturellement dans le cadre des graphes ;
elle donne aussi la description équivalente par la réunion $F = \bigcup F_i$,
dont les $F_i$ sont les composantes du graphe d'adjacence des faces, avec les
axiomes T1 et T3 du § VII.

\pagerange{136}{140}
\subsection*{Le « fait mirifique » (pages 136 à 140)}

But : un isomorphisme $Z \simeq (\mathbf{Z}/n\mathbf{Z})^{*}/\{\pm 1\}$. On
suppose (To 5) les sommets d'ordre fini $n \geq 3$. Par définition de
$A_\lambda$, l'étoile de $s$ pour $A_0$ est \emph{la même partie} de $S$ que
l'étoile de $\lambda s$ pour $A_\lambda$ ; seule change la structure de
polygone qu'y induit le graphe. Or une structure triangulaire qui satisfait
To 2 est déterminée par l'ensemble $\mathcal{E}_0$ de ses étoiles et, sur
chacune, la structure de $n$-gone induite ; toute arête est dans exactement
deux étoiles (une seule pour l'octaèdre). Il suffit donc de savoir comment
$A_\lambda$ change les polygones des étoiles. Et voici le « fait mirifique » :
\begin{quote}
  Soit $E$ un ensemble fini de cardinal $n \geq 3$. Le groupe $Z'_n =
  (\mathbf{Z}/n\mathbf{Z})^{*}/\{\pm 1\}$ opère canoniquement, et fidèlement,
  sur l'ensemble $\mathrm{Pol}(E)$ des structures de polygone sur $E$.
\end{quote}
Une structure de polygone \emph{orienté} est une permutation circulaire
$\omega$, ou encore une structure de torseur sous $\mathbf{Z}/n\mathbf{Z}$ (une
« droite affine sur $\Lambda_n$ ») ; $\mathrm{Aut}(\mathbf{Z}/n\mathbf{Z}) =
(\mathbf{Z}/n\mathbf{Z})^{*}$ y opère par $T_\nu \omega = \omega^\nu$ ;
$T_{-\nu}\omega = (T_\nu\omega)^{-1}$ est le polygone orienté opposé, et
l'action passe aux polygones non orientés en une action de $Z'_n$. Pour $n =
5$, l'élément non trivial change le pentagone en pentagramme\footnote{C'est
l'« étoilage » du dossier 78 (page 13, et pages 33 à 39, où le même groupe est
noté $\Psi_n$), sur les polygones réguliers du plan affine. Pour les cartes,
l'opération qui remplace en chaque sommet la rotation $\rho_0$ par
$\rho_0^\nu$ est connue aujourd'hui comme opérateur de Wilson, et les $\nu$ qui
donnent une carte isomorphe sont les \emph{exposants} de la carte (Wilson,
1979 ; Nedela et Škoviera, 1997). Cités de mémoire ; la page ne connaît pas
ces noms.}.

On pose alors (To 6) : pour tout $\lambda \in Z$ il existe $\nu \in Z'_n$ tel
que, pour tout sommet $s$, la structure de polygone induite par $A_\lambda$ sur
$\mathrm{Et}_{A_0}(s) = \mathrm{Et}_{A_\lambda}(\lambda s)$ soit $\alpha_{i_0,
s}^{\nu}$. Alors $A_\lambda = A_0^{(\nu)}$ est entièrement déterminé par $\nu$,
et (hors l'octaèdre) $\lambda$ aussi : $\lambda s$ est l'unique sommet dont
l'étoile pour $A_0^{(\nu)}$ est $\mathrm{Et}_0(s)$. D'où un homomorphisme
injectif $Z \hookrightarrow Z'_n$, dont l'image est formée de $\nu$ tels que
\[
  A_0^{(\nu)} = \coprod_{E \in \mathcal{E}_0} A_E^{(\nu)}
\]
soit encore l'ensemble des arêtes d'un 2-polyèdre --- « une propriété
remarquable, et qui pourrait être très forte, qui mérite d'être examinée pour
son propre mérite »\footnote{Dans le modèle (de nous) : pour $\lambda =
T_\mu$, l'étoile de $\dot x$ est $\{y + kx\}$ ; pour $A_0$, $y + kx$ est voisin
de $y + (k \pm 1)x$, pour $A_{T_\mu}$ de $y + (k \pm \mu)x$. Donc $\nu = \mu$,
et l'injection $Z \to Z'_n$ est l'isomorphisme $T_\mu \mapsto \mu$.}.

\pagerange{141}{145}
\subsection*{Les premiers cas, et le torseur des orientations (pages 141 à 145)}

La question n'a d'intérêt que si $Z'_n \neq 1$, c'est-à-dire $n \notin \{2, 3,
4, 6\}$. Pour $n = 5$, $Z'_5 \simeq \mathbf{Z}/2$ ; les seules triangulations
dont tous les sommets sont d'ordre $5$ sont l'icosaèdre et l'icosaèdre
« géminé », non orientable ; l'opération fait passer d'un icosaèdre à l'autre
dans le bi-icosaèdre. Mais pour l'icosaèdre géminé, le commutant de $G_0$ est
trivial, $Z = 1$ : la validité de l'opération $A_0 \mapsto A_0^{(\nu)}$ dépend
donc bien de l'axiomatique. Pour $n = 7$, $Z'_7 \simeq \mathbf{Z}/3$, engendré
par $\nu = 2$ : est-il vrai que pour \emph{tout} 2-polyèdre triangulé à sommets
d'ordre $7$, $A_0^{(2)}$ définisse encore un tel 2-polyèdre ? Même question pour
$n = 8$ et $n = 10$, où $Z'_n \simeq \mathbf{Z}/2$\footnote{La page dresse en
passant la liste des ordres de $Z'_n = \varphi(n)/2$. Deux valeurs sont à
corriger : $Z'_9$ est d'ordre $3$ (la page : $4$, et « $\mathbf{Z}/4$ ») ; $Z'_{15}$
est d'ordre $4$ (la page : $12$, lecture douteuse). Les autres sont justes :
$Z'_n = 1$ pour $n \in \{2, 3, 4, 6\}$, d'ordre $2$ pour $n \in \{5, 8, 10,
12\}$, d'ordre au moins $3$ pour $n \geq 13$ et $n \in \{7, 9, 11\}$.}.

Poursuivant l'axiomatique, la page pose (T 7) : $Z \to Z'_n$ est surjectif,
donc un isomorphisme ; et (T 8) : $P_0$ est régulier et orientable. On doit
alors pouvoir tout traduire en propriétés d'un groupe $G_0^{+}$ engendré par
$\rho_0, \rho_1, \rho_\infty$ avec $\rho_0^n = \rho_1^2 = \rho_\infty^3 = 1$,
$\rho_0$ d'ordre exactement $n$, sous-groupe d'indice $2$ d'un groupe $G_0$
engendré par des involutions $\tau_0, \tau_1, \tau_\infty$, avec $\rho_0 =
\tau_\infty\tau_1$, $\rho_1 = \tau_0\tau_\infty$, $\rho_\infty =
\tau_1\tau_0$\footnote{La relation $\rho_\infty\rho_1\rho_0 = 1$ découle alors
de ces égalités.}.

L'ensemble $\overrightarrow{\pi} = \coprod_{i \in \pi} \mu(i)$ des couples
(structure, orientation) devient un torseur sous $\Lambda_n^{*} =
(\mathbf{Z}/n\mathbf{Z})^{*}$ : pour $\omega \in \mu(i)$ et $\lambda \in
\Lambda_n^{*}$ d'image $\dot\lambda$, on pose $\lambda \cdot \omega = \omega'$,
où $\omega' \in \mu(\dot\lambda i)$ est l'orientation pour laquelle la
permutation circulaire $\omega'_E$ de chaque étoile est $\omega_E^{\lambda}$.
L'associativité est immédiate. L'indépendance par rapport à l'étoile choisie
demande que les « conventions d'orientation » des divers $P_i$ sur $G_0$
coïncident.

\pagerange{146}{150}
\subsection*{Les conventions d'orientation (pages 146 à 150)}

Un long aparté examine l'abélianisé de $G_0$ : comme $\dot\rho_\infty^3 = 1$
entraîne $\dot\tau_0 = \dot\tau_1$, il est engendré par $\dot\tau_1,
\dot\tau_\infty$ avec $(\dot\tau_1\dot\tau_\infty)^n = 1$ ; il est donc
d'ordre au plus $2$ si $n$ est impair --- un seul sous-groupe d'indice $2$ ---
et un quotient de $\mathbf{F}_2 \times \mathbf{F}_2$ si $n$ est pair, avec a
priori trois caractères que $G$ pourrait permuter. La page se demande si cela
arrive pour $n = 8$ (abélianisé de $\mathrm{GL}'(2, \mathbf{Z}/8)^{\pm}$), puis
pour $\widehat{\mathbf{Z}}$, et conclut par une réduction à $n = 4$ que le
caractère en question est trivial\footnote{Lecture mince : une grande partie
des pages 147 à 149 est barrée, et l'aparté entre doubles crochets se ferme
sur cette conclusion sans que le raisonnement soit lisible en entier.}.

L'argument direct qui suit est net (pages 149 et 150). $G_0^{+} =
\mathrm{Aut}^{+}(P_{i_0})$ est le seul sous-groupe d'indice $2$ de $G_0$ qui
contienne la rotation $\rho_0$ autour d'un sommet $s_0$ ; il suffit de voir que
$\rho_0 \in \mathrm{Aut}^{+}(P_{i_1})$, $i_1 = \lambda \cdot i_0$. Or $\rho_0$
préserve $E = \mathrm{Et}_0(s_0) = \mathrm{Et}_1(\lambda s_0)$ et y est une
rotation du polygone $A_0 \cap \mathfrak{P}_2(E)$, donc aussi du polygone
$A_1 \cap \mathfrak{P}_2(E) = (A_0 \cap \mathfrak{P}_2(E))^{(\lambda)}$ ; comme il
fixe $\lambda s_0$ (il fixe $s_0$ et commute à $\lambda$), c'est une rotation
de $P_{i_1}$ autour de $\lambda s_0$. Donc
\begin{quote}
  il existe un unique sous-groupe $G_0^{+}$ de $G_0$ tel que $G_0^{+} =
  \mathrm{Aut}^{+}(P_i)$ pour tout $i \in \pi$,
\end{quote}
ce qui règle aussi la question laissée ouverte page 120.

\pagerange{150}{152}
\subsection*{Le caractère cyclotomique (pages 150 à 152)}

$G = \mathrm{Aut}(S, F)$ opère sur $\overrightarrow{\pi}$ en commutant à la
structure de torseur sous $\Lambda_n^{*}$ ; il en résulte un homomorphisme
\[
  \chi : G \longrightarrow (\mathbf{Z}/n\mathbf{Z})^{*}, \qquad
  g(i, \omega) = \chi(g) \cdot (i, \omega) ,
\]
que la page appelle le « \emph{caractère cyclotomique} », et qui relève le
« caractère arithmétique » $\chi' : G \to Z'_n \simeq Z$. Son noyau est
$G_0^{+}$, d'où une suite exacte $1 \to G_0^{+} \to G \to
(\mathbf{Z}/n\mathbf{Z})^{*} \to 1$, où $Z \simeq Z'_n$ s'envoie dans
$(\mathbf{Z}/n\mathbf{Z})^{*}$ par $\dot\lambda \mapsto \lambda^2$, et $G_0 =
\chi^{-1}(\{\pm 1\})$ (pour l'octaèdre, il faut remplacer $Z$ par
$\{1\}$)\footnote{Dans le modèle, $\chi$ est le déterminant $\mathrm{GL}(M)'
\to (\mathbf{Z}/n\mathbf{Z})^{*}$, qui agit sur les orientations $e \in (\det
M)^{*}$. Le nom est celui de la page. Le rapprochement est de nous : sur les
points de $n$-torsion d'une courbe elliptique, l'accouplement de Weil identifie
le déterminant de l'action de Galois au caractère cyclotomique ; rien ici ne dit
que la page y pense.}.

On en lit $Z \cap G_0^{+} \simeq \bigl({}_2(\mathbf{Z}/n\mathbf{Z})^{*}\bigr)/\{\pm
1\}$, le sous-groupe des éléments d'ordre $2$ modulo $\pm 1$, isomorphe à
$\mathbf{F}_2^{\nu}$ avec
\[
  \nu = \nu(n) - 1 + \varepsilon, \qquad
  \varepsilon = 0 \text{ si } e_2 \in \{0, 2\}, \quad -1 \text{ si } e_2 = 1,
  \quad +1 \text{ si } e_2 \geq 3 ,
\]
où $\nu(n)$ est le nombre de facteurs premiers de $n$ et $e_2$ l'exposant de
$2$\footnote{Vérifié par nous : la $2$-torsion de $(\mathbf{Z}/p^e)^{*}$ est
$\{\pm 1\}$ pour $p$ impair, et celle de $(\mathbf{Z}/2^e)^{*}$ est triviale,
$\{\pm 1\}$, $\mathbf{F}_2^2$ pour $e = 1$, $2$, $\geq 3$.}. Ce groupe est
trivial exactement pour $n = p^e$, $n = 2p^e$ ($p$ premier impair) et $n = 4$ ;
pour $n = 2^e$, $e \geq 3$, il est d'ordre $2$, et la $2$-torsion de
$(\mathbf{Z}/n\mathbf{Z})^{*}$ est $\mathbf{F}_2^2$.

\pagerange{152}{156}
\subsection*{Le stabilisateur d'un sommet se scinde (pages 152 à 156)}

Les faces orientées et les arêtes orientées ont un type dans
$\overrightarrow{\pi}$, compatible à l'incidence (carrés cartésiens (67),
(68)). Soit $G_{s_0}$ le stabilisateur d'un sommet dans $G$ ; comme $G_0^{+}$
est transitif sur $S$, on a une extension
\[
  1 \longrightarrow G^{+}_{0 s_0} \longrightarrow G_{s_0}
  \xrightarrow{\ \chi\ } (\mathbf{Z}/n\mathbf{Z})^{*} \longrightarrow 1 ,
\]
où $G^{+}_{0s_0}$ est cyclique d'ordre $n$, non canoniquement. Soit
$\mathrm{Et}_F(s_0)$ l'ensemble des voisins de $s_0$ pour tous les types ;
$G_{s_0}$ y est transitif, le stabilisateur d'un voisin $t_0$ étant $\{1,
\tau_\infty\}$. $Z$ opère librement sur $S$ et préserve $\mathrm{Et}_F(s_0)$ ;
soit $D_0$ l'orbite de $t_0$ sous $Z$, et $T$ le stabilisateur de $D_0$ dans
$G_{s_0}$. Alors $T \cap G_0^{+} = \{1\}$ et $\chi$ induit un isomorphisme de
$T$ sur $(\mathbf{Z}/n\mathbf{Z})^{*}$, de sorte que
\[
  G_{s_0} \simeq G^{+}_{0s_0} \rtimes T
\]
est un produit semi-direct\footnote{La page 154 écrit « $Z \simeq
(\mathbf{Z}/n\mathbf{Z})^{*}$ » ; c'est $Z \simeq
(\mathbf{Z}/n\mathbf{Z})^{*}/\{\pm 1\}$, et c'est $T$ qui est isomorphe à
$(\mathbf{Z}/n\mathbf{Z})^{*}$. Dans le modèle (de nous), $G_{s_0}$ est
l'image des matrices $\bigl(\begin{smallmatrix} \pm 1 & b \\ 0 & d
\end{smallmatrix}\bigr)$, le groupe affine de la droite $\mathbf{Z}/n$, et $T$
celle des matrices diagonales $\mathrm{diag}(1, d)$ ; les gloses de la page 153,
« unipotents » et « Borel », vont dans ce sens.}. Si $g \in G^{+}_{0s_0} \cap
T$ envoie $t_0$ sur $\lambda t_0$, les types de $\{s_0, t_0\}$ et $\{s_0,
\lambda t_0\}$ sont égaux (car $g \in G_0$) et diffèrent de $\lambda$, donc
$\lambda = 1$ ; $g$ fixe alors $s_0$ et $t_0$, et c'est $1$ ou $\tau_\infty$,
mais $\chi(\tau_\infty) = -1$. La surjectivité vient de la transitivité sur les
voisins, corrigée au besoin par $\tau_\infty$. En passant (78) : si un élément
d'une $Z$-orbite $D$ est voisin de $s_0$, tous le sont, et $t \mapsto
\tau_A(\{s_0, t\})$ est une bijection de $D$ sur $\pi$.

\pagerange{156}{158}
\subsection*{Un groupe d'homothéties donné d'avance (pages 156 à 158)}

Cette décomposition est « formelle » à partir de T1 à T4 et de l'axiome de
régularité et d'orientabilité T8, sans les axiomes « arithmétiques » T5 à T7.
Dans cet esprit, la page repart d'un 2-polyèdre régulier $P_0$, de son groupe
$G_0$, et d'un groupe $Z \subset \mathrm{Centr}(G_0)$ \emph{donné d'avance},
vérifiant : pour tout $\lambda \in Z$ il existe $f$ centralisant $Z$ avec $\{s,
t\} \in A_0 \iff \{f^{-1}s, f^{-1}\lambda t\} \in A_0$. Par les
« considérations générales maintes fois répétées », $\pi = \{A_\lambda\}$ est un
torseur sous $Z$ (sauf pour l'octaèdre), le groupe $G$ des permutations qui
centralisent $Z$ et préservent $\pi$ opère par translations, et l'on a
l'extension
\[
  1 \longrightarrow G_0 \longrightarrow G \xrightarrow{\ \chi'\ } Z
  \longrightarrow 1 , \qquad \chi'|_Z : \lambda \mapsto \lambda^2 .
\]
Reste à définir l'équivalence de 2-polyèdres munis d'un groupe d'« homothéties ».

\pagerange{159}{164}
\section*{IX. Graphes à homothéties (pages 159 à 177)}

\subsection*{Définitions (pages 159 à 164)}

Un \emph{graphe à homothéties} est la donnée sur $S$ d'un graphe $A_0 \subset
\mathfrak{P}_2(S)$ et d'un sous-groupe $Z \subset \mathfrak{S}_S$ --- dont on
ne suppose pas qu'il respecte $A_0$ --- avec
\begin{itemize}
  \item (GH 1) pour tout $\lambda \in Z$ il existe $f \in \mathfrak{S}_S$
    centralisant $Z$ tel que $\{s, t\} \in A_0 \iff \{f^{-1}s, f^{-1}\lambda t\}
    \in A_0$ ;
\end{itemize}
c'est-à-dire : a) le graphe $A_\lambda$ ($\{s, t\} \in A_0 \iff \{s, \lambda t\}
\in A_\lambda$) est bien défini ; b) il est isomorphe à $A_0$ ; c) par un
isomorphisme qui commute à $Z$. On pose $\pi = \{A_\lambda \mid \lambda \in
Z\}$ ; $Z$ y opère par $T_\lambda$, avec $\{s, t\} \in A_i \iff \{s, \lambda t\}
\in A_{T_\lambda i}$. Deux graphes à homothéties de même $Z$ sont équivalents
si l'un est dans le $\pi$ de l'autre ; une classe est un \emph{multigraphe à
homothéties}. Ses automorphismes stricts (commutant à $Z$) forment $G$, ceux de
$A_0$ forment $G_0$, et l'on a
\[
  1 \longrightarrow G_0 \longrightarrow G \longrightarrow Z/\mathfrak{z}
  \longrightarrow 1 ,
\]
$\mathfrak{z}$ étant le noyau de $Z \to \mathfrak{S}_\pi$, formé des $\lambda$
qui ne changent aucune étoile --- trivial si tout sommet est déterminé par son
étoile. L'ensemble des étoiles ne dépend pas de la structure $A_i$ choisie.

Un graphe est \emph{régulier} si $G_0$ est transitif sur les arêtes orientées
(et le graphe connexe, sans sommet isolé). Alors $Z$ opère librement sur $S$,
les $A_i$ sont deux à deux disjoints (chacun est l'orbite d'une de ses arêtes
sous $G_0$), d'où une application type $A = \bigcup A_i \to \pi$. Le graphe est
« prépolyédral » si ses triangles font de lui un 2-polyèdre (toute arête dans
exactement deux triangles, étoiles connexes). Dans le cas régulier et
prépolyédral, $G_0$ a des générateurs canoniques $\tau_0, \tau_1, \tau_\infty$
relatifs à un repère $(s_0, t_0, u_0)$, $\rho_0$ est d'ordre $n$ (l'ordre des
sommets), et le stabilisateur de $s_0$ dans $G_0$ est le groupe diédral
$\widetilde{U} = \langle \tau_1, \tau_\infty \rangle \simeq \mathbf{D}_n$.

\pagerange{165}{174}
\subsection*{Traduction en théorie des groupes (pages 165 à 174)}

La page traduit en général la donnée (26) : a) un espace homogène $S$ sous
$G_0$ avec une origine $s_0$ ; b) un groupe $G \supset G_0$, $G_0$ distingué,
opérant sur $S$ ; c) un sous-groupe central $Z$ de $G$. La donnée a) est celle
du couple $(G_0, \widetilde{U} = G_{0s_0})$, avec $S \simeq G_0/\widetilde{U}$.
La donnée b) est celle de $B' = G_{s_0}$, avec $\widetilde{U} = G_0 \cap B'$ et
$B'/\widetilde{U} \simeq G/G_0 = Z'$ (le carré $\widetilde{U} \subset G_0$,
$B' \subset G$ est cartésien et cocartésien) ; on préfère la décrire par
l'extension $1 \to \widetilde{U} \to B' \to Z' \to 1$, l'inclusion $\widetilde{U}
\subset G_0$ et une action de $B'$ sur $G_0$ prolongeant l'action intérieure de
$\widetilde{U}$ --- des groupes « petits et bien compris », $Z'$ étant
$(\mathbf{Z}/n\mathbf{Z})^{*}$ ou son quotient par $\pm 1$, $\widetilde{U}$
cyclique ou diédral\footnote{Autrement dit (de nous), $G$ est le quotient du
produit semi-direct $G_0 \rtimes B'$ par l'image antidiagonale de
$\widetilde{U}$.}. Pour c), le commutant de $G_0$ dans $\mathfrak{S}_S$ étant
$N_{G_0}(\widetilde{U})/\widetilde{U}$, un sous-groupe $Z$ du commutant
correspond à un sous-groupe $\widetilde{B}''$ avec $\widetilde{U} \triangleleft
\widetilde{B}'' \subset N_{G_0}(\widetilde{U})$, $Z = \widetilde{B}''/\widetilde{U}$
--- en termes de $S$, $\widetilde{B}'' = \{g \in G_0 \mid g s_0 \in Z s_0\}$ ---
et que $Z$ commute à tout $G$ se lit : $B'$ normalise $\widetilde{B}''$ et opère
trivialement sur $Z$.

Le résumé (50) : a) $\widetilde{U} \triangleleft \widetilde{B}'' \subset G_0$, $Z =
\widetilde{B}''/\widetilde{U}$ ; b) un homomorphisme $Z \to Z'$ et une extension
$B'$ de $Z'$ par $\widetilde{U}$ dont $\widetilde{B}''$ soit l'image
réciproque, $\widetilde{B}'' \simeq B' \times_{Z'} Z$ ; c) une action de $B'$ sur
$G_0$ dont la composée avec $\psi : \widetilde{B}'' \to B'$ soit l'action
intérieure. On récupère $G$, $S \simeq G_0/\widetilde{U} \simeq G/B'$, et $Z$
comme l'image de $\varphi(g) = g\,\psi(g)^{-1} = \psi(g)^{-1} g$, qui est un
homomorphisme trivial sur $\widetilde{U}$ --- la page en vérifie la
multiplicativité en utilisant que $Z$ est commutatif\footnote{La vérification
(page 173) est écrite « pour $g, g' \in B'$ » ; elle porte sur $g, g' \in
\widetilde{B}''$. Une note marginale de la page 172 ajoute des conditions
« oubliées » sur l'action $\theta$.}.

\pagerange{174}{177}
\subsection*{Le cas des 2-polyèdres réguliers (pages 174 à 177)}

Dans le cas qui intéresse le dossier : $G_0$ est un quotient du groupe
cartographique triangulé
\[
  \mathfrak{G}_{\infty,3} = \langle \tau_0, \tau_1, \tau_\infty \mid \tau_0^2 =
  \tau_1^2 = \tau_\infty^2 = (\tau_1\tau_0)^3 = (\tau_0\tau_\infty)^2 = 1
  \rangle
\]
(relation d'ordre $3$ sur $\tau_1\tau_0 = \rho_\infty$\footnote{La page 175 écrit $(\tau_1\tau_\infty)^3$ sous l'accolade
  « $\rho_\infty$ » ; avec $\rho_\infty = \tau_1\tau_0$ (page 164, (21)), la
  relation d'ordre $3$ porte sur $\tau_1\tau_0$. Écrite comme sur la page,
  elle ferait de la rotation $\rho_0$ autour d'un sommet un élément d'ordre
  $3$.}) ;
$\widetilde{U} = \langle \tau_1, \tau_\infty \rangle \simeq \mathbf{D}_n$ ; $B'$
est un sous-groupe de $\mathrm{Aut}(G_0 \supset \widetilde{B}'' \supset
\widetilde{U})$ contenant l'image de $\widetilde{B}''$ et opérant trivialement
sur $Z$ ; et un isomorphisme $Z' \simeq Z$ par lequel $Z \to Z'$ devient $\lambda
\mapsto \lambda^2$ : l'extension $\widetilde{B}''$ est l'image réciproque de
$B'$ par l'élévation au carré. Reste à exprimer l'axiome GH 1 : la condition ne
dépend que de l'image de $f$ dans $G/G_0$, on peut prendre $f \in B'$ avec
$\dot f = \alpha(\lambda)$, et l'équivalence $\{s, t\} \in A_0 \iff \{f^{-1}s,
f^{-1}\lambda t\} \in A_0$ se réécrit, avec $s = us_0$, $t = vs_0$, $t_0 = \rho
s_0$, comme une condition sur des classes modulo $\widetilde{U}$ et $B'$
(62). « Donc la question à vérifier s'écrit » : la phrase reste en suspens, et
la page suivante ouvre le § X.

\pagerange{178}{189}
\section*{X. Antipodismes généralisés, et présentation autoduale (pages 178 à 221)}

\subsection*{1. Points antipodiques et homothéties (pages 178 à 189)}

Soit $S$ un espace homogène sous un groupe $G$, et $Z$ le commutant de $G$ dans
$\mathfrak{S}_S$. Deux points sont \emph{$G$-équivalents} s'ils ont le même
stabilisateur ; $Z$ opère librement, et les classes sont les $Z$-orbites. Avec
une origine $s_0$, $S \simeq G/U$, $U = G_{s_0}$, et $Z \simeq N_G(U)/U$ ; plus
généralement, les applications de $S$ dans lui-même qui commutent à $G$
correspondent aux points fixes $S^{U}$, c'est-à-dire aux classes $aU$ avec
$\mathrm{int}(a)U \supset U$, par $f \mapsto f(s_0)$, la translation à droite
renversant l'ordre des produits (14). On a $Z s_0 = S^{U}$ quand $S$ est fini,
quand les stabilisateurs sont finis, ou quand $G$ est un groupe de Lie réel à
stabilisateurs fermés n'ayant qu'un nombre fini de composantes. La page appelle
\emph{antipodiques} à $s_0$ les points de $S^{U}$, \emph{strictement
antipodiques} les points $G$-équivalents, et antipodismes les applications qui
commutent à $G$.

Appliquons cela aux sommets d'un 2-polyèdre connexe régulier, dont les sommets
sont d'ordre fini\footnote{La page 182 écrit « dans le cas où les ordres des
sommets sont pairs, … i.e. dans le cas (16) b » ; le cas (16 b) est celui des
stabilisateurs finis, et il faut lire « finis ».} : antipodismes et antipodismes
stricts coïncident, et $Z$ est le groupe des antipodismes. Un $\lambda \in Z$ est
une \emph{homothétie} si
\[
  \text{(H1)} \qquad \{s, t\} \in A_0 \iff \{\lambda s, \lambda^{-1} t\} \in A_0 ,
  \qquad \text{c'est-à-dire} \qquad
  \lambda\,\mathrm{Et}(t) = \mathrm{Et}(\lambda^{-1} t)
\]
(« J'ignore si cette condition est toujours satisfaite ! »). Il suffit de
le vérifier pour une arête orientée $(s_0, t_0)$ ; avec $t_0 = \rho s_0$, $\rho
\in G_0$, $\lambda s_0 = a s_0$, $a \in N_{G_0}(\widetilde{U})$, la condition
s'écrit
\[
  \rho(a)\,a \in \widetilde{U}_1 \widetilde{U}_0 , \qquad
  \widetilde{U}_0 = G_{0s_0}, \ \widetilde{U}_1 = G_{0t_0} = \rho(\widetilde{U}),
\]
où $\rho(a) = \rho a \rho^{-1}$, et de même dans $G_0^{+}$ avec les
stabilisateurs directs (24').

L'homothétie est \emph{stricte} si de plus $A_\lambda$ est isomorphe à $A_0$.
Les étoiles de $A_0$ et $A_\lambda$ sont alors les mêmes parties de $S$, avec
des centres différents : le centre pour $A_\lambda$ est l'antipodique, par
$\lambda$, du centre pour $A_0$. Le sous-groupe des rotations autour de $s$ est
l'unique sous-groupe cyclique d'ordre $n$ de $G_{0s} \simeq \mathbf{D}_n$, le
même que pour $\lambda s$ ; les deux structures de polygone de l'étoile ont donc
le même groupe de rotations, et se déduisent l'une de l'autre par un élément
$\nu \in Z'_n$, indépendant du sommet : $A_\lambda = \bigcup_E A_{E,0}^{(\nu)}$
(29). C'est le « fait mirifique » de la page 138, retrouvé par un autre chemin.

Pour qu'une telle torsion donne un graphe triangulé, il faut a) que toute arête
de $A_\lambda$ soit dans une étoile, b) que chaque étoile soit un polygone pour
$A_\lambda$. La condition a) revient à ce que $\lambda^{-1}s$ soit à distance au
plus $3$ de $s$ ; en fait la distance vaut exactement $3$, sauf si $\lambda$ est
l'identité ou l'antipodie de l'octaèdre (distance $2$) : elle ne peut valoir $1$
(page 110, (31)), et la valeur $2$ voudrait dire que $s$ et $\lambda^{-1}s$ ont
la même étoile, ce qui n'arrive que pour l'octaèdre\footnote{C'est le lemme de
la page 275, énoncé plus tard : deux sommets distincts de même étoile forment
une double pyramide, qui n'est régulière que pour l'octaèdre. Le passage de
« un voisin commun » à « même étoile » utilise que $G_{0s}$ fixe aussi
$\lambda^{-1}s$ et est transitif sur l'étoile de $s$.}. L'examen de b) est
commencé (page 189) et laissé là.

\pagerange{190}{197}
\subsection*{2. Les contours invariants d'une étoile (pages 190 à 197)}

Soit $(S, A_0)$ un graphe où toute arête est dans exactement deux triangles, les
étoiles étant des contours d'ordre $\geq 3$, et $G_0$ un groupe
d'automorphismes transitif sur les arêtes orientées, simplement transitif sur
les repères triangulaires. Le graphe est connu par l'ensemble $\Sigma$ de ses
étoiles et les contours qu'il y induit ; la page cherche toutes les structures
de graphe $A_i$ de même groupe et de mêmes étoiles. Par transitivité, une telle
structure est connue par le contour qu'elle induit sur \emph{une} étoile $E_0$,
invariant par le stabilisateur de $E_0$. L'image de ce stabilisateur dans
$\mathfrak{S}_{E_0}$ est le groupe diédral $D_{E_0} \simeq \mathbf{D}_n$ du
polygone, de sous-groupe cyclique $\Lambda_{E_0}$ d'ordre $n$, qui ne dépend
que de $(S, G_0, \Sigma)$.

\emph{Proposition} (page 195). Soit $E$ un torseur sous un groupe cyclique
$\Lambda_E \subset \mathfrak{S}_E$ d'ordre $n$. Les structures de contour
orienté sur $E$ invariantes par $\Lambda_E$ correspondent aux éléments $u \in
\Lambda_E$, $u$ étant la rotation du contour : une structure de contour orienté
est une permutation $u$, et elle est invariante si $u$ commute à $\Lambda_E$,
dont le commutant, pour une action simplement transitive d'un groupe
commutatif, est lui-même.

\emph{Corollaires.} Les polygones orientés invariants ($u$ générateur) forment
un torseur sous $(\mathbf{Z}/n\mathbf{Z})^{*}$ ; les contours non orientés
invariants s'identifient à $(\Lambda_E \smallsetminus \{0, \sigma\})/\pm 1$, et
les polygones invariants forment un torseur sous $Z'_n$\footnote{Ici
« contour » s'entend au sens large : pour $u$ d'ordre $d < n$, c'est une réunion
de $n/d$ cycles de longueur $d$. On exclut les cas dégénérés $u = 0$ et $u =
\sigma$, l'élément d'ordre $2$ quand $n$ est pair, qui ne donnent que des points
ou des paires ; une note marginale juge cette exclusion « inutilement
restrictive », et le monoïde multiplicatif de $\mathbf{Z}/n\mathbf{Z}$ opère
sur l'ensemble complet.}. En termes d'un générateur $\omega_0$, la structure
associée à $\nu$ est définie par la paire $\{\omega_0^{\nu}, \omega_0^{-\nu}\}$.

Les systèmes de contours $(C_E)_{E \in \Sigma}$ invariants par $G_0$
correspondent donc aux éléments de $(\Lambda_{E_0} \smallsetminus \{0,
\sigma\})/\pm 1$, et pour chacun $A = \bigcup_E A_E$ est stable par $G_0$,
transitif sur ses arêtes orientées. Reste à savoir lequel provient d'une
structure du type cherché.

\pagerange{198}{201}
\subsection*{3. Un graphe comme incidence entre sommets et étoiles (pages 198 à 201)}

Un graphe strict est une relation symétrique $R_0 \subset S \times S$ sans
point diagonal ; l'étoile de $s$ est $R_0(s)$. Soit $S^{\vee}$ l'ensemble des
étoiles (« cosommets ») et $R \subset S \times S^{\vee}$ la relation
d'incidence. Si chaque sommet est déterminé par son étoile --- le graphe est
« séparé » ---, l'application $\varphi_0 : s \mapsto R_0(s)$ est une bijection
$S \to S^{\vee}$, et l'on reconstruit le graphe par $R_0 = (\mathrm{id} \times
\varphi_0)^{-1}(R)$. Il en résulte une équivalence de groupoïdes entre graphes
séparés et systèmes $(S, S^{\vee}, R, \varphi_0)$ où les deux applications $S
\to \mathfrak{P}(S^{\vee})$, $S^{\vee} \to \mathfrak{P}(S)$ sont injectives et
où $\varphi_0$ est une bijection « symétrique » ($(s, \varphi_0 t) \in R \iff (t,
\varphi_0 s) \in R$) et « antiréflexive » ($(s, \varphi_0 s) \notin R$). La
description fait apparaître une \emph{dualité} : $(S, S^{\vee}, R, \varphi_0)
\mapsto (S^{\vee}, S, {}^{t}R, \varphi_0^{-1})$, le graphe dual étant
canoniquement isomorphe au graphe de départ par $\varphi_0$\footnote{Dans le
langage des géométries d'incidence (rapprochement de nous), $\varphi_0$ est une
\emph{polarité} de la structure d'incidence $(S, S^{\vee}, R)$ sans point
absolu, et le graphe est le graphe de cette polarité.}.

\pagerange{201}{206}
\subsection*{4. Prégraphes et jauges ; l'exemple linéaire (pages 201 à 206)}

Le « jeu nouveau » est de partir d'une incidence $(S, S^{\vee}, R)$ sans
$\varphi_0$ --- un \emph{prégraphe} --- et d'étudier \emph{toutes} les
bijections $\varphi : S \to S^{\vee}$ symétriques et antiréflexives, les
\emph{jauges}, dont $\varphi_0$ n'est qu'une « origine ». Écrivant $\varphi =
\varphi_0 \circ \lambda$, $\lambda \in \mathfrak{S}_S$, la condition devient
\[
  \{s, t\} \in A_0 \iff \{\lambda^{-1}s, \lambda t\} \in A_0 , \qquad
  \{s, \lambda s\} \notin A_0 ,
\]
c'est-à-dire $\lambda\,\mathrm{Et}(s) = \mathrm{Et}(\lambda^{-1}s)$ : on retrouve
exactement les homothéties.

L'exemple, en « digression » (pages 203 à 206) : $M$ et $M^{\vee}$ libres de
rang $2$ en dualité, $S = M^{*}/\{\pm 1\}$, $S^{\vee} = M^{\vee *}/\{\pm 1\}$,
$R = \{(\dot x, \dot{\check x}) \mid \langle x, \check x \rangle = \pm 1\}$. Une
jauge $e$ de $M$ donne la forme alternée $\psi_e$, $x \wedge y = \psi_e(x, y)\,
e$, et l'isomorphisme $\Phi_e : M \to M^{\vee}$, $x \wedge y = \langle x,
\Phi_e(y)\rangle\, e$ ; l'étoile de $s$ dans le graphe $A_e$ est $\{t \mid (s,
\varphi_e(t)) \in R\}$. Ce graphe est séparé sauf dans le « cas octaédral
généralisé » $2\Lambda = \{0, 2\} \neq \{0\}$ (pour $\Lambda$ semi-local ; pour
$\mathbf{Z}/n\mathbf{Z}$, sauf $n = 4$), et il est alors décrit par $(S,
S^{\vee}, R, \varphi_e)$. Le groupe $Z'_\Lambda$ opère librement sur $S$ et
$S^{\vee}$, et $(s, \check s) \in R \iff (\lambda s, \lambda^{-1}\check s) \in R$.
La bijection $\varphi_e$ ne dépend que de $\{\pm e\}$, et comme $\Phi_{\lambda e}
= \lambda^{-1}\Phi_e$, il vaut mieux indexer les jauges par $(\det
M^{\vee})^{*}/\{\pm 1\}$, ce qui donne $\varphi_{\lambda\alpha}(s) =
\lambda\,\varphi_\alpha(s) = \varphi_\alpha(\lambda s)$ (55).

\pagerange{207}{213}
\subsection*{Homothéties d'un prégraphe ; un repentir (pages 207 à 213)}

Une homothétie d'un prégraphe est un couple $(\lambda_S, \lambda_{S^{\vee}})$
tel que $(s, \check s) \in R \iff (\lambda s, \lambda^{-1}\check s) \in R$ ;
quand $R$ est séparé des deux côtés, chacune des deux permutations détermine
l'autre. L'ensemble des homothéties est un sous-ensemble « un peu vicieux » de
$\mathfrak{S}_S \times \mathfrak{S}_{S^{\vee}}$, qui n'est pas un sous-groupe :
si $R$ est le graphe d'une bijection, c'est l'antidiagonale $\{(u, u^{-1})\}$. On
change donc de normalisation : on prend les automorphismes $(\lambda,
\check\lambda)$ du prégraphe, et la jauge $\varphi = \check\lambda\varphi_0 =
\varphi_0\lambda^{-1}$ --- dans l'exemple, les scalaires opèrent sur $M^{\vee}$
par contragrédience. Deux homothéties d'un graphe séparé composent en une
homothétie si elles commutent (66) ; on se donne donc un sous-groupe
\emph{commutatif} $Z$ d'automorphismes, et l'on considère l'ensemble $\Pi_Z$
des jauges compatibles, $\check\lambda\varphi\lambda = \varphi$ pour tout
$(\lambda, \check\lambda) \in Z$, sur lequel $Z$ opère librement.

La page affirme alors que $\Pi_Z$ est un torseur sous le centralisateur
$\mathcal{G}$ de $Z$, conclut « Ouf », puis entoure le mot « torseur » et écrit
« faux »\footnote{Ce qui reste établi : $Z$ opère librement sur $\Pi_Z$.
L'application $g \mapsto g_{S^{\vee}}\varphi_0 = \varphi_0 g_S^{-1}$ est
injective sur l'ensemble des $g$ pour lesquels $\varphi_0$ est une
« $g$-jauge », mais cet ensemble n'est pas un groupe. La Scholie de la page 213,
qui répète l'énoncé, est marquée « caduc » dans la marge. Un « Lemme » écrit en oblique dans la marge de la page 212 énumère des conditions sur un tel $g$ ; il n'est lisible que par fragments, et on ne le restitue pas.}. Les multigraphes à
homothéties envisagés auparavant sont les sous-$Z$-torseurs de $\Pi_Z$.

\pagerange{214}{221}
\subsection*{Boucles, torsion, transport de structure (pages 214 à 221)}

Le graphe tordu $A_\lambda$ peut avoir des boucles ($t = \lambda s$ avec $\{s,
t\} \in A_0$). Exemple : $S^{\vee} = S$, $R$ la diagonale ; les jauges sont les
involutions de $S$, les automorphismes les $(\lambda, \lambda)$, et $\Pi_Z$ est
formé des involutions qui normalisent $Z$ en y induisant $\lambda \mapsto
\lambda^{-1}$. Si $S$ est un torseur sous $Z$ (noté additivement), ce sont les
$\varphi_\lambda(z) = \lambda - z$, qui ont un point fixe si et seulement si
$\lambda \in 2Z$ : selon l'origine choisie, on a ou non une « vraie »
structure de graphe.

Pour un sous-$Z$-torseur $\Pi$ de $\Pi_Z$, l'évaluation $(\varphi, s) \mapsto
\varphi(s)$ donne une bijection de $Z$-ensembles
\[
  \Pi \wedge^{Z} (S!) \xrightarrow{\ \sim\ } S^{\vee} ,
\]
où $S!$ est $S$ muni de l'action inverse ; elle commute à $g \in \mathcal{G}$
si et seulement si $\varphi$ commute à $g$, et le plus grand groupe pour lequel
c'est le cas est $\mathcal{G}_0 = \{g \in \mathcal{G} \mid \varphi g_S =
g_{S^{\vee}}\varphi\}$ --- dans l'exemple linéaire, $\mathrm{GL}^{\pm}(M)$
\footnote{Une note marginale le dit. Vérification (de nous) : avec $\check u =
{}^{t}u^{-1}$, on a $\check u\,\Phi_e = \det(u)^{-1}\,\Phi_e\, u$, de sorte que
$\Phi_e$ commute à $u$ modulo $\pm 1$ si et seulement si $\det u = \pm 1$.}. La
structure complète $(S, S^{\vee}, R, Z, \Pi)$ est autoduale : $(S^{\vee}, S,
{}^{t}R, {}^{t}Z, \Pi^{-1})$ vérifie les mêmes conditions, et $S \simeq \Pi
\wedge^{Z} (S^{\vee}!)$. Par transport de structure, un automorphisme $g$ envoie
$\varphi$ sur $g(\varphi) = g_{S^{\vee}}\varphi g_S^{-1}$ et $A_\varphi$ sur
$A_{g(\varphi)}$ ; pour $g \in Z$, $g(\varphi) = T_{g^2}\varphi$ --- le carré, une
fois de plus.

La page 220, « Retour sur la condition d'isomorphie », reprend avec $\mathcal{A}
= \mathrm{Aut}(\Gamma)$ : les permutations $f$ telles que $f(A_{\varphi_0})$ soit
de la forme $A_\varphi$, $\varphi \in \Pi$, forment un sous-groupe $G$ du
centralisateur de $Z$, contenant $\mathcal{G}_0 = \mathrm{Aut}(S, A_\varphi)$ et
$Z$, avec $1 \to \mathcal{G}_0 \to G \xrightarrow{\delta} Z$, $\delta(\lambda)
= \lambda^2$ sur $Z$. La page 221, barrée, est une première rédaction de ce qui
suit.

\pagerange{222}{227}
\section*{XI. Récapitulation (pages 222 à 270)}

\subsection*{1. Prégraphes, jauges, multigraphes (pages 222 à 227)}

Soit $\Gamma = (S, S^{\vee}, R)$ un prégraphe (« sommets », « cosommets »,
« incidence »). Deux applications $\varphi, \psi : S \to S^{\vee}$ sont
\emph{associées} si $(s, \psi t) \in R \iff (t, \varphi s) \in R$ ; une
\emph{jauge} est une $\varphi$ associée à elle-même, \emph{stricte} si elle est
bijective. Les jauges strictes correspondent aux isomorphismes $(\varphi,
\varphi^{-1})$ de $\Gamma$ sur son dual $(S^{\vee}, S, {}^{t}R)$. Une jauge
$\varphi$ définit le graphe $A_\varphi = \{\{s, t\} \mid (s, \varphi t) \in
R\}$, boucles permises ; si $\Gamma$ est séparé en $S^{\vee}$, on obtient des
équivalences de groupoïdes : graphes et prégraphes séparés en $S^{\vee}$ munis
d'une jauge surjective ; graphes séparés et prégraphes biséparés munis d'une
jauge stricte. Le groupe $\mathcal{A} = \mathrm{Aut}(\Gamma)$ s'identifie alors
aux permutations de $S$ qui envoient étoile sur étoile.

Les homothéties du graphe $A_\varphi$ correspondent aux $\lambda \in \mathcal{A}$
avec $\lambda_{S^{\vee}}\varphi = \varphi\lambda_S^{-1}$ (« $\lambda$-jauges »)
(21) ; une jauge bijective compatible à $\lambda$ et à $\mu$ l'est à $\lambda\mu$
si et seulement si $\lambda\mu = \mu\lambda$, de sorte qu'un groupe qui admet une
jauge stricte compatible est commutatif (23). Pour $Z \subset \mathcal{A}$
commutatif, $Z$ opère librement sur l'ensemble $\Pi_Z$ des $Z$-jauges strictes
par $T_\lambda\varphi = \lambda_{S^{\vee}}\varphi = \varphi\lambda_S^{-1}$, ce
qui tord le graphe en $T_\lambda A_\varphi = \{\{s, \lambda t\}\}$\footnote{La
page 226, (29), affirme encore que $g \mapsto g_{S^{\vee}}\varphi_0$ est une
\emph{bijection} de l'ensemble $\mathcal{G}_{\varphi_0}$ des $g$ du
centralisateur pour lesquels $\varphi_0$ est une $g$-jauge sur $\Pi_Z$.
L'injectivité est claire ; la surjectivité est ce que la page 212 avait
barré d'un « faux », et rien ici ne l'établit.}. Un « multigraphe » à groupe
d'homothéties $Z$ est un sous-$Z$-torseur $\pi \subset \Pi_Z$ ; on a $S^{\vee}
\simeq \pi \wedge^{Z} (S!)$, et dualement.

\pagerange{228}{234}
\subsection*{Rayons (pages 228 à 234)}

Les orbites de $Z$ dans $S$ sont les \emph{rayons}, celles dans $S^{\vee}$ les
\emph{corayons} ; toutes les $\varphi \in \pi$ induisent la même bijection
$\varphi^\flat : S/Z \to S^{\vee}/Z$. On suppose $Z$ libre sur $S$ et
$S^{\vee}$ : les rayons sont des $Z$-torseurs. Les conditions équivalentes
(41), (41'), (41''), (44), (48) disent :
\begin{itemize}
  \item une étoile et un rayon se rencontrent en au plus un point ;
  \item étant donné une jauge, $A_0 \cap A_\lambda = \emptyset$ pour $\lambda
    \neq 1$, de sorte que les $A_i$ sont disjoints et qu'on a un type $\tau_A :
    A \to \pi$ ;
  \item pour deux rayons $D$, $D'$ incidents, l'adjacence est une bijection $D
    \simeq D'$ qui \emph{renverse} les structures de torseur, d'où $D
    \wedge^{Z} D' \simeq \mathbf{1}$ ; et le type vérifie $\tau(\lambda s, \mu s')
    = \lambda\mu\,\tau(s, s')$, d'où un isomorphisme de $Z$-torseurs $D \wedge D'
    \simeq \pi$.
\end{itemize}
Sur une feuille rapportée, numérotée « 119 bis » (pages 233 et 234), la page en
tire une description globale : un graphe à homothéties, avec $Z$ libre et cette
condition, équivaut à a) un graphe $\gamma^\flat$ sur l'ensemble $S^\flat$ des
rayons, b) un $Z$-torseur $S$ au-dessus de $S^\flat$, c) pour chaque arête
$\{s^\flat, s'^\flat\}$ un isomorphisme $p_1^{*}S \wedge p_2^{*}S \simeq
\mathbf{1}$ ; un multigraphe, à la même chose avec un $Z$-torseur absolu $\pi$
et des isomorphismes $p_1^{*}S \wedge p_2^{*}S \simeq \pi$\footnote{En termes
d'aujourd'hui (de nous) : un torseur sous un groupe abélien au-dessus des sommets
d'un graphe, avec des identifications le long des arêtes, est la donnée d'un
« graphe de voltages », et le graphe à homothéties en est le revêtement
régulier associé.}. Exemple : pour $M$ libre de rang $2$, $S^\flat = M^{*}/\Lambda^{*}
\subset \mathbf{P}(M)(\Lambda)$ (égalité si $\operatorname{Pic}\Lambda = 0$),
$\gamma^\flat$ relie deux points de la droite projective distincts en tout point
de $\operatorname{Spec}\Lambda$ ($x \wedge y$ inversible), le torseur a pour
fibre $\Lambda^{*}/\{\pm 1\}$, et l'isomorphisme c) vient de $\det M \simeq D
\otimes D'$ pour $M = D \oplus D'$\footnote{Pour $\Lambda = \mathbf{Z}$,
$\gamma^\flat$ est le graphe de Farey sur $\mathbf{P}^1(\mathbf{Q})$ (remarque de
nous). Le tableau de la page 272 compte les points de $\mathbf{P}(M)$ pour $n =
2$ à $5$.}.

\pagerange{235}{237}
\subsection*{2. La condition d'isomorphie (pages 235 à 237)}

Soit $\mathcal{M} = (\Gamma, Z, \pi)$ un « prégraphe multijaugé ». Ses
automorphismes qui centralisent $Z$ forment le groupe $G$ des $g$ du
centralisateur avec $g(\pi) = \pi$, qui opère sur $\pi$ par translations :
\begin{tikzcd}
1 \arrow[r] & G_0 \arrow[r] & G \arrow[r, "\delta"] & Z \\
 & & Z \arrow[u, hook] \arrow[ur, "\lambda \mapsto \lambda^{2}"'] &
\end{tikzcd}
où $G_0 = \operatorname{Ker}\delta$ s'identifie aux automorphismes du graphe
$A_{\varphi_0}$ qui commutent à $Z$, et où $\delta(\lambda) = \lambda^2$ parce
que $\lambda(\varphi) = \lambda_{S^{\vee}}\varphi\lambda_S^{-1} =
\lambda_{S^{\vee}}^2\varphi$. Sont équivalents : a) $\delta$ est surjectif ; b)
$G$ est transitif sur $\pi$ ; c) chaque $A_\lambda$ est isomorphe à $A_0$ par
une permutation commutant à $Z$ ; d) l'axiome GH 1. C'est la \emph{condition
d'isomorphie}\footnote{La page 237 dit, en achevant la phrase, que « la condition
d'homogénéité est satisfaite » ; le titre du numéro, et la page 221, appellent
cette condition « d'isomorphie », et le numéro suivant réserve le mot
« homogénéité » à une autre condition. On suit le titre.}.

\pagerange{238}{247}
\subsection*{3. Homogénéité, et traduction en théorie des groupes (pages 238 à 269)}

$\mathcal{M}$ est \emph{homogène} si $G_0$ est transitif sur $R$ --- ce qui
revient à la transitivité sur les arêtes orientées de $A_{\varphi_0}$ --- et sur
$S$ et $S^{\vee}$. Alors $Z$ opère librement (si $\Gamma$ est séparé), et la
condition des rayons (41) est satisfaite : si $A_0 \cap A_\lambda \neq
\emptyset$, la transitivité donne $A_0 = A_\lambda$, d'où $\lambda = 1$.

Pour tout décrire en théorie des groupes, la page part d'une situation
« épinglée » : un graphe $\gamma_0 = (S, A_0)$, un groupe $H$ d'automorphismes
transitif sur les arêtes orientées, une arête orientée origine $(s_0, t_0)$, et
un groupe commutatif $Z$ d'homothéties commutant à $H$ ; $G$ est alors le groupe
des permutations qui normalisent $H$, centralisent $Z$ et envoient $A_0$ sur un
$A_\lambda$, et $H$ y est distingué (64). Les données $(H, Z, (s_0, t_0))$
déterminent $A = \{\{u s_0, u\lambda t_0\}\}$, sous les conditions : a) $H$
transitif sur $S$ et commutant à $Z$ ; b) pour tout $\lambda \in Z$, il existe $u
\in H$ avec $u(s_0, t_0) = (\lambda s_0, \lambda^{-1}t_0)$\footnote{La page 242
écrit $u(s_0, t_0) = (s_0, \lambda t_0)$, avec un signe surchargé devant
$s_0$. Ainsi écrite, la condition contredirait la condition de fidélité (72)
qui suit ; la remarque de la même page sur le cas $s_0 = t_0$ (« $\lambda^2 =
1$ ») et la définition de $T_H$ page 245 demandent $(\lambda s_0,
\lambda^{-1}t_0)$, qui est la condition d'homothétie (H1).} ; et l'on suppose
que $Z$ opère fidèlement sur $\pi$ (72).

Soit $U_0 = H_{s_0}$, $U_1 = H_{t_0}$, $\mathfrak{z}_H = U_0 \cap U_1$, et $\sigma
\in H$ échangeant $s_0$ et $t_0$ : $U_1 = \sigma U_0 \sigma^{-1}$, $\sigma^2 \in
\mathfrak{z}_H$. Un graphe homogène muni d'une arête orientée équivaut à la
donnée de $(H, U_0, \sigma)$ avec ces relations et $\bigcap_u uU_0u^{-1} = 1$ :
$S = H/U_0$, $t_0 = \sigma U_0$, les arêtes orientées $= H/\mathfrak{z}_H$ ---
$\sigma$ n'étant défini que modulo $\mathfrak{z}_H$. Quand $H$ est
\emph{simplement} transitif sur les arêtes orientées, $\mathfrak{z}_H = 1$,
$\sigma$ est unique et $\sigma^2 = 1$\footnote{C'est la description
d'aujourd'hui des graphes arc-transitifs comme graphes de classes (Sabidussi),
et, pour $\mathfrak{z}_H = 1$, des graphes arc-réguliers. Noms de nous.}.

La donnée de $Z$ se lit alors sur le sous-groupe
\[
  T_H = \{u \in H \mid \exists \lambda \in Z,\ us_0 = \lambda s_0,\ ut_0 =
  \lambda^{-1}t_0\} ,
\]
avec $1 \to \mathfrak{z}_H \to T_H \to Z \to 1$ : $T_H$ normalise $U_0$ et
$U_1$, $T_H \cap U_0 = T_H \cap U_1 = \mathfrak{z}_H$, $T_H U_0$ est le
stabilisateur du rayon $Z s_0$, $T_H U_0 \cap T_H U_1 = T_H$, et $\sigma$
normalise $T_H$ en y induisant $t \mapsto t^{-1}$ modulo $\mathfrak{z}_H$. La
situation est « particulièrement jolie » quand $\mathfrak{z}_H = 1$ (89) : la
donnée de $Z$ équivaut à celle d'un sous-groupe $T_H$ de $H$ tel que a) $\sigma$
normalise $T_H$ et y induit l'inversion, b) $T_H$ normalise $U_0$, c)
l'application $U_0 \times T_H \times U_1 \to H$, $(u_0, t, u_1) \mapsto u_0 t
u_1$, est injective\footnote{Dans le modèle (de nous), $H = \mathrm{SL}(2,
\mathbf{Z}/n)/\{\pm 1\}$, $s_0 = \dot e_1$, $t_0 = \dot e_2$ : $U_0$ est le
groupe unipotent supérieur, $U_1$ l'inférieur, $T_H$ le tore diagonal $\simeq
(\mathbf{Z}/n)^{*}/\{\pm 1\}$, $\sigma = \bigl(\begin{smallmatrix} 0 & -1 \\ 1 &
0 \end{smallmatrix}\bigr)$ ; c) est l'unicité de la décomposition de Gauss (ici
sous la forme $UDL$) sur la grosse cellule de Bruhat.}.

\pagerange{248}{260}
\subsection*{Le groupe $G$ reconstruit (pages 248 à 260)}

Dans ce cas, la page reconstruit tout le diagramme (64). $T_G$, stabilisateur
dans $G$ des deux rayons $D_0 = Zs_0$, $D_1 = Zt_0$, est une extension de $Z
\times Z$ par $\mathfrak{z} = G_{(s_0, t_0)}$ ; il contient les sous-groupes
$T_0$ (qui fixe $s_0$) et $T_1$ (qui fixe $t_0$)\footnote{La page 249 définit
$T_0$ par « $gs_0 = s_0$, $gt_0 \in Zt_0$ » et $T_1$ symétriquement, tout en les
appelant $\operatorname{Ker}\delta_0$, $\operatorname{Ker}\delta_1$ ; les deux
étiquettes sont échangées par rapport aux définitions, qu'on suit.}, avec $T_G =
T_0 T_1 = T_0 Z = Z T_1$, $T_0 \cap T_1 = T_0 \cap Z = \mathfrak{z}$, et $\sigma$
échange $T_0$ et $T_1$. On a $T_G \cap H = T_H$, $T_G$ centralise $T_H$,
normalise $U_0$ et $U_1$, et
\[
  G_{D_0} = T_G U_0, \quad G_{s_0} = T_0 U_0, \quad G_0 = \mathfrak{z} \cdot H,
  \quad G = T_0 \cdot H = T_1 \cdot H
\]
(produits semi-directs)\footnote{La page 252 écrit « $G = T_0.H = T'_0.H$ » ; le
second facteur, lu avec doute, est $T_1$, par symétrie.}, l'application $U_0
\times T_G \times U_1 \to G$ restant injective. Avec $\widetilde{Z} = G/H$, on a
$0 \to \mathfrak{z} \to \widetilde{Z} \xrightarrow{c} Z$, et $c$ est surjectif
si et seulement si la condition d'isomorphie est satisfaite ; un diagramme $3
\times 3$ d'inclusions, à carrés cartésiens et cocartésiens, résume le tout
(113).

$T_G$ est une extension $1 \to T_H \simeq Z \to T_G \to \widetilde{Z} \to 1$,
scindée par $T_0$ et par $T_1$ (sections $u_0$, $u_1$), avec $u_1(\tilde\lambda)
= u_0(\tilde\lambda)\, u(c(\tilde\lambda))$ ; elle est commutative. Si
$\widetilde{Z}$ est commutatif, $T_G \simeq (T_0 \times T_1)/\mathfrak{z}$,
$\mathfrak{z}$ plongé antidiagonalement, et $\delta$ se lit par $c$. L'inclusion
de $Z$ dans $T_G$ se décrit par l'homomorphisme $c' : Z \to \widetilde{Z}$,
$c \circ c' = (\lambda \mapsto \lambda^2)$, par $\lambda = u(\lambda)\,
u_0(c'(\lambda)) = u(\lambda^{-1})\, u_1(c'(\lambda))$ ; on voudrait voir $Z$
comme l'image de la diagonale de $\widetilde{Z} \times \widetilde{Z}$, ce qui
demande $z^2 = 1$ sur $\mathfrak{z}$ (131) et $c'(c(\tilde\lambda)) =
\tilde\lambda^2$ (133) ; en général on n'a que $c(\tilde\lambda) \equiv
u_0(\tilde\lambda)\, u_1(\tilde\lambda)$ modulo $\mathfrak{z}$\footnote{Dans le
modèle (de nous) : $G = \mathrm{GL}(2, \mathbf{Z}/n)/\{\pm 1\}$, $\widetilde{Z} =
(\mathbf{Z}/n)^{*}$ par le déterminant, $Z = Z'_n$, $c$ la projection,
$\mathfrak{z} = \{1, \mathrm{diag}(1, -1)\}$, $c'(\lambda) = \lambda^2$,
$u_0(d) = \mathrm{diag}(1, d)$, $u_1(d) = \mathrm{diag}(d, 1)$, et
$u_0(d)u_1(d)$ est la classe du scalaire $d$ : (131), (132) et (133) y sont
vrais.}.

\pagerange{261}{269}
\subsection*{Par les automorphismes de $H$ ; le groupe $N_H$ (pages 261 à 269)}

Reste à tout exprimer en termes des seules données $(H, U_0, T_H, \sigma)$ :
$\sigma^2 = 1$, $U_0 \cap U_1 = 1$, $\sigma$ inverse $T_H$, $T_H$ normalise $U_0$,
$U_0 \times T_H \times U_1 \to H$ injectif. Un élément de $T_0$ est connu par
l'automorphisme $\bar g$ de $H$ qu'il induit ; $T_0$ s'identifie au sous-groupe
$\widetilde{Z}$ de $\mathrm{Aut}(H)$ formé des automorphismes qui normalisent
$U_0$, centralisent $T_H$, et envoient $\sigma$ dans $\sigma T_H$ ; on a alors
$\bar g(\sigma) = \sigma\, u(\delta(g))$, ce qui décrit $c$, et la conjugaison par
$t \in T_H$ donne $c'_0 : T_H \to \widetilde{Z}$ avec $c_0(c'_0(t)) = t^{-2}$,
puisque $t\sigma t^{-1} = \sigma t^{-2}$. Le sous-groupe $\mathfrak{z}$
correspond aux automorphismes qui normalisent $U_0$ et centralisent $T_H$ et
$\sigma$.

La page introduit enfin $N_H = T_H \cup \sigma T_H$, qui contient $T_H$ avec
l'indice $2$ (hors le cas dégénéré $\sigma \in T_H$, où $U_0 = U_1 = 1$ et $T_H$
est un $2$-groupe abélien élémentaire), et énonce une équivalence (163) : un
multigraphe homogène non dégénéré, avec un sommet marqué et un rayon incident,
équivaut à un système $(H, U_0, T_H, N_H)$ tel que a) $T_H$ normalise $U_0$ et
$T_H \cap U_0 = 1$ (d'où $B_H = T_H U_0$) ; b) $T_H$ est d'indice $2$ dans
$N_H$, et les $\sigma(U_0)$, $\sigma \in N_H \smallsetminus T_H$, sont un même
$U_1$ ; c) tout $\sigma \in N_H \smallsetminus T_H$ est d'ordre $2$ ; d) $U_0
\cap U_1 = 1$. « On aimerait vérifier aussi la formule familière » $B_H \cap N_H =
T_H$, qui équivaut à $\sigma \notin B_H$ ; la tentative est barrée\footnote{Les
lettres $B$, $N$, $T$, $U$ et la « formule familière » sont celles des systèmes
de Tits (BN-paires) : $B \cap N = T$, $N/T$ d'ordre $2$ en rang $1$. Le
rapprochement est de nous ; la page ne nomme pas Tits, et sa note marginale, qui
semble dire que $(U_0, T_H, N_H)$ se déduisent de $(B_H, N_H)$, est d'une
lecture très incertaine.}.

\pagerange{270}{270}
\subsection*{4. Cas des graphes triangulants (page 270)}

Un graphe sans boucles est \emph{triangulant} si toute arête est contenue dans
un triangle ; un multigraphe l'est si chacune de ses structures l'est, et il
suffit que l'une le soit sous la condition d'isomorphie, « mais il n'est pas dit
qu'il en soit ainsi dans tous les cas ». Le texte suivi du dossier s'arrête là,
à la page 138 de l'auteur.

\pagerange{272}{273}
\section*{Feuilles détachées (pages 272 à 295)}

\subsection*{Un tableau (pages 272 et 273)}

Une fiche, déchirée en bas, dresse pour $n = 2$ à $5$ les invariants du § I
(on note $[X]$ le cardinal de $X$)\footnote{La page écrit $\mathrm{GP}(1,
\mathbf{Z}/n\mathbf{Z})$ pour $\mathrm{PGL}(2, \mathbf{Z}/n\mathbf{Z})$, comme le
titre du dossier 82, « $\mathrm{SO}(3) \simeq \mathrm{GP}(1)$ ». Les deux points
d'interrogation de la page se lèvent : $\mathrm{PGL}(2, \mathbf{Z}/4) =
\mathrm{GL}(M)'$ (les scalaires inversibles y sont $\pm 1$) est, par le théorème
de la page 89, le groupe de l'octaèdre, $\mathfrak{S}_4 \times \{\pm 1\}$ ; et
$\mathrm{PGL}(2, \mathbf{F}_5) \simeq \mathfrak{S}_5$ est classique. La ligne
$[F_M^{*}]$ compte les faces de tous les types (le bi-icosaèdre en a $40$),
$[F_M^{\circ}]$ celles d'un seul type ; ses chiffres sont surchargés. La page écrit $\mathfrak{S}_4^{+}$ et $\mathfrak{S}_5^{+}$ pour les groupes alternés.} :
\[
\begin{array}{l|c|c|c|c}
  n & 2 & 3 & 4 & 5 \\ \hline
  [\mathrm{SL}(2, \mathbf{Z}/n)] & 6 & 24 & 48 & 120 \\
  {}[\mathrm{GL}(2, \mathbf{Z}/n)] & 6 & 48 & 96 & 480 \\
  {}[\mathrm{SL}(2, \mathbf{Z}/n)'] & 6 & 12 & 24 & 60 \\
  \mathrm{SL}(2, \mathbf{Z}/n)' & \mathfrak{S}_3 & \mathfrak{A}_4 &
    \mathfrak{S}_4 & \mathfrak{A}_5 \\
  {}[\mathrm{PGL}(2, \mathbf{Z}/n)] & 6 & 24 & 48 & 120 \\
  \mathrm{PGL}(2, \mathbf{Z}/n) & \mathfrak{S}_3 & \mathfrak{S}_4 &
    \mathfrak{S}_4 \times \{\pm 1\} & \mathfrak{S}_5 \\
  {}[M^{*}] & 3 & 8 & 12 & 24 \\
  {}[S_M] & 3 & 4 & 6 & 12 \\
  \text{type} & \text{trièdre} & \text{tétraèdre} & \text{octaèdre} &
    \text{icosaèdre} \\
  {}[\mathbf{P}(M)] & 3 & 4 & 6 & 6 \\
  {}[F_M^{*}] & 1 & 4 & 8 & 40 \\
  {}[F_M^{\circ}] & 1 & 4 & 8 & 20
\end{array}
\]
Pour $n = 2$, où $2 = 0$, le « trièdre » a trois sommets et une seule face. Au
verso, la seule inscription « $(\mathbf{Z}/n\mathbf{Z})^{2*}$ ».

\pagerange{275}{280}
\subsection*{Un sommet est-il déterminé par son étoile ? (pages 275 à 280)}

1) \emph{Combinatoire.} Soit $P = (S, F)$ un 2-polyèdre triangulé
connexe\footnote{La page 275 écrit « convexe » ; l'argument utilise la
connexité.}, et $s \neq s'$ deux sommets de même étoile $E$. Alors $S = \{s, s'\}
\cup E$ et $P$ est le double cône sur le polygone $E$ : chaque côté de $E$ borde
déjà une face de sommet $s$ et une face de sommet $s'$, de sorte que ces faces
forment une surface fermée, qui est tout $P$. Les pointes sont d'ordre $n =
\operatorname{card} E$, les autres sommets d'ordre $4$ : si tous les sommets ont
le même ordre, $n = 4$ et $P$ est un octaèdre, où deux sommets ont la même
étoile si et seulement s'ils sont opposés. Dans tous les cas, l'involution $a$
qui échange les sommets de même étoile (« qui mérite le nom
d'antipodisme ? ») est un automorphisme.

2) \emph{Algèbre.} Soit $\Lambda$ semi-local, $u, u' \in M^{*}$ tels que, pour
tout $x \in M^{*}$, $u \wedge x$ inversible entraîne $u' \wedge x \in \{\pm u
\wedge x\}$. Alors $u' = \pm u$, sauf si
\[
  2\Lambda = \{0, 2\} \neq \{0\} ,
\]
c'est-à-dire si $\Lambda$ est une $\mathbf{Z}/4$-algèbre, extension de
l'$\mathbf{F}_2$-algèbre $\Lambda/2\Lambda$ par l'idéal $2\Lambda$, de carré
nul, de dimension $1$ sur $\mathbf{F}_2$\footnote{La page parle d'une
« connexion » $\chi : \Lambda/2\Lambda \to \mathbf{F}_2$ : c'est le caractère par
lequel $\Lambda/2\Lambda$ opère sur la droite $2\Lambda \simeq \mathbf{F}_2$, dont
le noyau est l'idéal maximal annulateur de $2$.}. La démonstration : la
première moitié de l'hypothèse dit que $u$ et $u'$ sont colinéaires dans chaque
corps résiduel ; on trouve une base $(e, f)$ avec $u = e + f$, $u' = e \pm f$ ;
le cas $u' = e - f$ force le corps résiduel à être de caractéristique $2$ et
$2\alpha$ ou $2(1 - \alpha)$ nul pour tout $\alpha$, d'où $2\Lambda \subset \{0,
2\}$. Dans le cas exceptionnel, $u' = u + \xi$ avec $\xi$ l'un des deux éléments
de $2M \smallsetminus \{0, 2u\}$, et les deux choix sont opposés : la classe de
$u'$ dans $S_M$ est déterminée par $u$. On obtient une involution $a$ de $S_M$,
sans point fixe, avec $\mathrm{Et}_i(a(u)) = \mathrm{Et}_i(u)$ pour tout type
$i$, et $a \in \mathrm{GL}'(M)^{\pm}$, de déterminant $-1$\footnote{Pour
$\Lambda = \mathbf{Z}/4$, on vérifie (de nous) que $a$ est la classe de la
matrice $\bigl(\begin{smallmatrix} 1 & 2 \\ 2 & -1 \end{smallmatrix}\bigr)$, de
déterminant $-5 = -1$.}. Pour $\Lambda = \mathbf{Z}/4$ c'est l'antipodie de
l'octaèdre, « retrouvée par voie d'algèbre linéaire » : c'est l'exception du §
VI, et celle que les §§ VII et VIII rencontraient sous le nom d'octaèdre. La
construction marche dès que $2\Lambda = \{0, 2\} \neq \{0\}$, mais rien ne dit
alors qu'il n'y a pas d'autres sommets de même étoile.

La page 280, au crayon, essaie des conditions de connexité (deux sommets reliés
par une chaîne de longueur $4$) et note : « Sommet déterminé par son étoile ?
Pas vrai pour $\Lambda = \mathbf{Z}/4$ (octaèdre) ».

\pagerange{282}{295}
\subsection*{Brouillons : $\lambda$ et $\lambda^2$ (pages 282 à 295)}

Les dernières feuilles sont des calculs au crayon, sans pagination, où les
mêmes formules reviennent ; on n'en restitue que ce qu'elles établissent
lisiblement\footnote{Lecture mince. Une phrase isolée, à l'encre bleue, le
long du bord de la page 282 --- « irréelle … au beau milieu de délire
mathématico-philosophique » --- n'a pas de lien visible avec les calculs.}.
L'étoile de $s = \dot x$ est un torseur sous $\Lambda$ par $y \dotplus \zeta = y +
\zeta x$ ; l'étoile de $\lambda^{-1}s$ l'est par $y' \dotplus \zeta = y' +
\zeta\lambda^{-1}x$. La multiplication par $\lambda$ envoie la première sur la
seconde, avec
\[
  \lambda(y \dotplus \zeta) = \lambda y \dotplus \lambda^2\zeta ,
\]
et la page souligne : « on ne voit ici que $\lambda^2$, pas $\lambda$ ! ». Mais la
relation d'incidence entre les deux étoiles, $y \wedge y' \in \{\pm e\}$, s'écrit,
pour des origines $y'_0 = \lambda y_0$,
\[
  \lambda\zeta - \lambda^{-1}\zeta' \in \{\pm 1\} , \qquad \text{soit} \qquad
  \zeta' = \lambda^2\zeta \pm \lambda ,
\]
de sorte que les écarts $\zeta' - \lambda^2\zeta$ forment l'ensemble $\{\pm
\lambda\}$ : la multiplication elle-même se voit sur les incidences, au signe
près. C'est la réponse, à l'état de calcul, à l'obstacle de la page 109
(« $\chi_s(\dot\lambda) = \lambda^2$ ») ; les feuilles comptent encore les
arêtes et les étoiles ($N/2$ et $N/n$), notent que $\mathrm{diag}(\lambda^{-1},
\lambda)$ normalise les groupes unipotents et y agit par $\lambda^2$, et
s'arrêtent sur une définition de $\nu(A_0, \lambda)$ laissée en suspens. Le
dossier finit là.

\end{document}
